% Generated mechanically from frozen input inputs/p12/ja/moduli.tex.
% Do not edit this generated file; edit the bound locale source instead.

% OK, start here.
%

\setcounter{chapter}{107}
\chapter{モジュライスタック}


\phantomsection
\label{moduli-section-phantom}


% BEGIN EXPANDED MODULE ch108_01.tex
\section{序論}
\label{moduli-section-introduction}

\noindent
本章では，モジュライ空間およびモジュライスタックの基本的性質を検証する．
対象となるものは
$\mathit{Hom}$，$\mathit{Isom}$，$\Cohstack_{X/B}$，
$\Quotfunctor_{\mathcal{F}/X/B}$，$\Hilbfunctor_{X/B}$，
$\Picardstack_{X/B}$，$\Picardfunctor_{X/B}$，$\mathit{Mor}_B(Z, X)$，
$\Spacesstack'_{fp, flat, proper}$，$\Polarizedstack$，および
$\Complexesstack_{X/B}$ などである．
% P12-ERR-0190: Japanese supplies the predicate omitted by the frozen English sentence.
これらが代数空間または代数スタックであることは，適切な仮定の下で既に示した；
Quot の次の節を参照せよ：
\ref{quot-section-hom}，
\ref{quot-section-isom}，
\ref{quot-section-stack-coherent-sheaves}，
\ref{quot-section-not-flat}，
\ref{quot-section-quot}，
\ref{quot-section-hilb}，
\ref{quot-section-picard-stack}，
\ref{quot-section-picard-functor}，
\ref{quot-section-relative-morphisms}，
\ref{quot-section-stack-of-spaces}，
\ref{quot-section-polarized}，および
\ref{quot-section-moduli-complexes}．
$\textit{Curves}$ と記す曲線のスタックは Quot の節
\ref{quot-section-curves} で導入したものであり，曲線のモジュライを扱う章で論じる；
Moduli of Curves, Section \ref{moduli-curves-section-stack-curves} を参照せよ．

\medskip\noindent
ある意味で，本章は Grothendieck の講義
\cite{Gr-I}，
\cite{Gr-II}，
\cite{Gr-III}，
\cite{Gr-IV}，
\cite{Gr-V}，および
\cite{Gr-VI} の足跡に従うものである．







\section{規約と言葉の濫用}
\label{moduli-section-conventions}

\noindent
Properties of Stacks, Section
\ref{stacks-properties-section-conventions} で導入した規約と言葉の濫用を，
引き続き用いる．特に断らない限り，基底スキームは $\Spec(\mathbf{Z})$ とする．






\section{Hom と Isom の性質}
\label{moduli-section-hom-isom}

\noindent
$f : X \to B$ を有限表示な代数空間の射とする．
$\mathcal{F}$ と $\mathcal{G}$ は準連接な $\mathcal{O}_X$-加群と仮定する．
$\mathcal{G}$ が有限表示で，$B$ 上平坦であり，その台が $B$ 上 proper ならば，
$\mathit{Hom}(\mathcal{F}, \mathcal{G})$ で定める関手
$$
T/B \longmapsto \Hom_{\mathcal{O}_{X_T}}(\mathcal{F}_T, \mathcal{G}_T)
$$
は $B$ 上アフィンな代数空間である．$\mathcal{F}$ も有限表示ならば，
$\mathit{Hom}(\mathcal{F}, \mathcal{G}) \to B$ は有限表示である．
Quot, Proposition \ref{quot-proposition-hom} を参照せよ．

\medskip\noindent
$\mathcal{F}$ と $\mathcal{G}$ の双方が有限表示で，$B$ 上平坦，かつ台が $B$ 上 proper ならば，
部分関手
$$
\mathit{Isom}(\mathcal{F}, \mathcal{G}) \subset
\mathit{Hom}(\mathcal{F}, \mathcal{G})
$$
は $B$ 上有限表示のアフィンな代数空間である．
Quot, Proposition \ref{quot-proposition-isom} を参照せよ．




\section{連接層のスタックの性質}
\label{moduli-section-stack-coherent-sheaves}

\noindent
$f : X \to B$ を分離的かつ有限表示な代数空間の射とする．このとき
$\Cohstack_{X/B}$ は，proper な台をもつ連接加群の平坦族をパラメータ付けるスタックであり，
代数的である．Quot, Theorem \ref{quot-theorem-coherent-algebraic-general} を参照せよ．

\begin{lemma}
\label{moduli-lemma-coherent-diagonal-affine-fp}
$\Cohstack_{X/B}$ の $B$ 上の対角射はアフィンであり，有限表示である．
\end{lemma}

\begin{proof}
対角射が代数空間によって表現可能であることは Quot, Lemma
\ref{quot-lemma-coherent-diagonal} で示されている．
その証明から，次を示せばよい：
$\mathit{Isom}(\mathcal{F}, \mathcal{G}) \to T$
がアフィンかつ有限表示であることを，有限表示な
$\mathcal{O}_{X_T}$-加群の組
$\mathcal{F}$，$\mathcal{G}$ で，$T$ 上平坦かつ台が $T$ 上 proper なものについて示せばよい．
これは Section \ref{moduli-section-hom-isom} で扱った．
\end{proof}

\begin{lemma}
\label{moduli-lemma-coherent-qs-lfp}
射 $\Cohstack_{X/B} \to B$ は準分離的で，局所有限表示である．
\end{lemma}

\begin{proof}
$\Cohstack_{X/B} \to B$ が準分離的であることを確認するには，
その対角射が準コンパクトかつ準分離的であることを示せばよい．
これは Lemma \ref{moduli-lemma-coherent-diagonal-affine-fp} から直ちに従う．
$\Cohstack_{X/B} \to B$ が局所有限表示であることを示すには，
$\Cohstack_{X/B} \to B$ が極限を保つことを示せばよい；Limits of Stacks, Proposition
\ref{stacks-limits-proposition-characterize-locally-finite-presentation} を参照せよ．
これは Quot, Lemma \ref{quot-lemma-coherent-limits} から従う（細部は省略する）．
\end{proof}

\begin{lemma}
\label{moduli-lemma-coherent-existence-part}
$X \to B$ が proper かつ有限表示であると仮定する．このとき
$\Cohstack_{X/B} \to B$ は，付値判定法（Morphisms of Stacks, Definition
\ref{stacks-morphisms-definition-existence}）の存在部分を満たす．
\end{lemma}

\begin{proof}
基底変換を行えば，これは直ちに次の問題に帰着する：付値環 $R$ とその分数体 $K$，
代数空間 $X$ で $R$ 上 proper なもの，および連接な
$\mathcal{O}_{X_K}$-加群 $\mathcal{F}_K$ が与えられたとき，有限表示な
$\mathcal{O}_X$-加群 $\mathcal{F}$ で，$R$ 上平坦かつ一般ファイバーが $\mathcal{F}_K$ であるものが存在することを示せ．
Flatness on Spaces, Theorem \ref{spaces-flat-theorem-finite-type-flat} により，
有限型準連接 $\mathcal{O}_X$-加群 $\mathcal{F}$ が $R$ 上平坦ならば有限表示であることに注意する．
% P12-ERR-0191: Japanese uses a grammatical let/notation construction.
一般ファイバーの埋め込みを $j : X_K \to X$ と書く．
これはアフィン射 $\Spec(K) \to \Spec(R)$ の基底変換なので，$j$ はアフィンである．
したがって $j_*\mathcal{F}_K$ は準連接である．
$$
j_*\mathcal{F}_K = \colim \mathcal{F}_i
$$
と，その有限型準連接 $\mathcal{O}_X$-部分加群によるフィルター付き余極限を書ける；
Limits of Spaces, Lemma \ref{spaces-limits-lemma-directed-colimit-finite-type} を参照せよ．
$j_*\mathcal{F}_K$ は $K$-ベクトル空間の $X$ 上の層なので，$\Spec(R)$ 上平坦である．
ゆえに各 $\mathcal{F}_i$ は $R$ 上平坦である．実際，付値環上の平坦性は torsion-free であることと同値であり
(More on Algebra, Lemma \ref{more-algebra-lemma-valuation-ring-torsion-free-flat})，
torsion-free 性は部分加群へ遺伝する．
最後に，写像 $j^*\mathcal{F}_i \to \mathcal{F}_K$ がある $i$ について同型であることを示す必要がある．
$j^*j_*\mathcal{F}_K = \mathcal{F}_K$（細部は省略）であり，$j^*$ は完全なので，
写像 $j^*\mathcal{F}_i \to \mathcal{F}_K$ はすべての $i$ について単射である．
$j^*$ は余極限と可換するので，
$\mathcal{F}_K = j^*j_*\mathcal{F}_K = \colim j^*\mathcal{F}_i$ である．
$\mathcal{F}_K$ は連接（すなわち有限表示）だから，ある $i$ が存在して，$j^*\mathcal{F}_i$ は
アフィン \'etale 被覆上の (有限個の)生成元をすべて含む．この被覆は $X$ の被覆である．
したがって，$j^*\mathcal{F}_i \to \mathcal{F}_K$ は $i$ が十分大きいとき全射である．
\end{proof}

\begin{lemma}
\label{moduli-lemma-coherent-functorial}
$B$ を代数空間とする．$\pi : X \to Y$ を，$B$ 上分離的かつ有限表示な代数空間の間の準有限射とする．
このとき $\pi_*$ は射
$\Cohstack_{X/B} \to \Cohstack_{Y/B}$ を誘導する．
\end{lemma}

\begin{proof}
$(T \to B, \mathcal{F})$ を $\Cohstack_{X/B}$ の対象とする．
次を主張する：
\begin{enumerate}
\item[(a)] $(T \to B, \pi_{T, *}\mathcal{F})$ は $\Cohstack_{Y/B}$ の対象であり，
\item[(b)] $T' \to T$ に対して
$\pi_{T', *}(X_{T'} \to X_T)^*\mathcal{F} =
(Y_{T'} \to Y_T)^*\pi_{T, *}\mathcal{F}$ である．
\end{enumerate}
(b) により，この構成が求める関手
$\Cohstack_{X/B} \to \Cohstack_{Y/B}$ を定めることが保証される．

\medskip\noindent
$i : Z \to X_T$ を，$\mathcal{F}$ の zeroth fitting ideal が切り出す閉部分空間とする
(Divisors on Spaces, Section \ref{spaces-divisors-section-fitting-ideals})．
% P12-ERR-0192: The frozen base B is preserved rather than silently changed to T.
仮定により $Z \to B$ は proper である（Derived Categories of Spaces, Section
\ref{spaces-perfect-section-proper-over-base} を参照）．
一方，$i$ は有限表示である
(Divisors on Spaces, Lemma
\ref{spaces-divisors-lemma-fitting-ideal-of-finitely-presented} および
Morphisms of Spaces, Lemma
\ref{spaces-morphisms-lemma-closed-immersion-finite-presentation})．
準連接 $\mathcal{O}_Z$-加群 $\mathcal{G}$ で，有限型かつ
$i_*\mathcal{G} = \mathcal{F}$ となるものが存在する
(Divisors on Spaces, Lemma
\ref{spaces-divisors-lemma-on-subscheme-cut-out-by-Fit-0})．
実際，Descent on Spaces, Lemma
\ref{spaces-descent-lemma-finite-finitely-presented-module} により，
$\mathcal{G}$ は $\mathcal{O}_Z$-加群として有限表示である．
% P12-ERR-0193: The frozen base B is preserved rather than silently changed to T.
$\mathcal{G}$ は $B$ 上平坦であることに注意する；例えば $\mathcal{G}$ と $\mathcal{F}$ の茎が一致するからである
(Morphisms of Spaces, Lemma \ref{spaces-morphisms-lemma-stalk-push-closed})．
$\pi_T \circ i : Z \to Y_T$ は準有限射の合成なので準有限であり，
% P12-ERR-0194: The frozen unmatched closing parenthesis is preserved.
$\pi_{T, *}\mathcal{F} = (\pi_T \circ i)_*\mathcal{G})$ であることにも注意する．
$i$ はアフィンなので，$i_*$ の形成は基底変換と可換する
(Cohomology of Spaces, Lemma \ref{spaces-cohomology-lemma-affine-base-change})．
したがって $B$ を $T$ に，$X$ を $Z$ に，$\mathcal{F}$ を $\mathcal{G}$ に，$Y$ を $Y_T$ に置き換えてよく，
次の段落で扱う場合に帰着する．

\medskip\noindent
$X \to B$ が proper であると仮定する．すると $\pi$ は
Morphisms of Spaces, Lemma \ref{spaces-morphisms-lemma-universally-closed-permanence} により proper であり，
したがって More on Morphisms of Spaces, Lemma
\ref{spaces-more-morphisms-lemma-characterize-finite} により有限である．
有限射はアフィンなので，(b) は Cohomology of Spaces, Lemma
\ref{spaces-cohomology-lemma-affine-base-change} により成り立つ．
一方，Morphisms of Spaces, Lemma
\ref{spaces-morphisms-lemma-finite-presentation-permanence} により $\pi$ は有限表示である．
ゆえに Descent on Spaces, Lemma
\ref{spaces-descent-lemma-finite-finitely-presented-module} により
$\pi_{T, *}\mathcal{F}$ は有限表示である．
最後に，例えば
Cohomology of Spaces, Lemma \ref{spaces-cohomology-lemma-stalk-push-finite} を用いて茎を調べれば，
$\pi_{T, *}\mathcal{F}$ は $B$ 上平坦である．
\end{proof}

\begin{lemma}
\label{moduli-lemma-coherent-open}
$B$ を代数空間とする．$\pi : X \to Y$ を，$B$ 上分離的かつ有限表示な
代数空間の間の開埋め込みとする．このとき Lemma
\ref{moduli-lemma-coherent-functorial} の射
$\Cohstack_{X/B} \to \Cohstack_{Y/B}$ は開埋め込みである．
\end{lemma}

\begin{proof}
省略する．ヒント：$\mathcal{F}$ が $T$ 上の $\Cohstack_{Y/B}$ の対象で，
$t \in T$ に対して $\text{Supp}(\mathcal{F}_t) \subset |X_t|$ ならば，
$t$ のある近傍に属する $t' \in T$ についても同じことが成り立つ．
\end{proof}

\begin{lemma}
\label{moduli-lemma-coherent-closed}
$B$ を代数空間とする．$\pi : X \to Y$ を，$B$ 上分離的かつ有限表示な
代数空間の間の閉埋め込みとする．このとき Lemma
\ref{moduli-lemma-coherent-functorial} の射
$\Cohstack_{X/B} \to \Cohstack_{Y/B}$ は閉埋め込みである．
\end{lemma}

\begin{proof}
$\mathcal{I} \subset \mathcal{O}_Y$ を，$X$ を $Y$ の閉部分空間として切り出すイデアル層とする．
$\pi_*$ は，準連接 $\mathcal{O}_X$-加群の圏と，$\mathcal{I}$ によって零化される
準連接 $\mathcal{O}_Y$-加群の圏との間の同値を誘導することを想起せよ；
Morphisms of Spaces, Lemma \ref{spaces-morphisms-lemma-i-star-equivalence} を参照せよ．
% P12-ERR-0195: Japanese supplies the word “change” omitted by the frozen English phrase.
$T \to B$ による基底変換の後にも，同じことが mutatis mutandis に成り立つ．
ただし $\mathcal{I}$ はイデアル層
$\mathcal{I}_T = \Im((Y_T \to Y)^*\mathcal{I} \to \mathcal{O}_{Y_T})$ に置き換える．
Lemma \ref{moduli-lemma-coherent-functorial} の証明を解析すると，
$\Cohstack_{X/B} \to \Cohstack_{Y/B}$ の本質的像は，ちょうど
$\mathcal{F}$ が $\mathcal{I}_T$ によって零化されるような対象
$\xi = (T \to B, \mathcal{F})$ からなることが分かる．
言い換えれば，$\xi$ が本質的像に属することと，乗法写像
$$
\mathcal{F} \otimes_{\mathcal{O}_{Y_T}} (Y_T \to Y)^*\mathcal{I}
\longrightarrow
\mathcal{F}
$$
が零であり，さらに任意の基底変換 $T' \to T$ の後にも同様であることとは同値である．
次に注意せよ：
$$
(Y_{T'} \to Y_T)^*(
\mathcal{F} \otimes_{\mathcal{O}_{Y_T}} (Y_T \to Y)^*\mathcal{I}) =
(Y_{T'} \to Y_T)^*\mathcal{F} \otimes_{\mathcal{O}_{Y_{T'}}}
(Y_{T'} \to Y)^*\mathcal{I})
$$
したがって，$T'$ 上で乗法写像が零となる条件は，Flatness on Spaces, Lemma
\ref{spaces-flat-lemma-F-zero-closed-proper} により，$T$ の閉部分空間によって表現可能である．
\end{proof}

\begin{situation}[数値的不変量]
\label{moduli-situation-numerical}
$f : X \to B$ は本節の序文におけるものとする．$I$ を集合とし，各 $i \in I$ に対して
$E_i \in D(\mathcal{O}_X)$ を perfect とする．
$\Cohstack_{X/B}$ の対象 $(T \to B, \mathcal{F})$ が与えられたとき，
$E_i$ の $X_T$ への導来引き戻しを $E_{i, T}$ と記す．
$D(\mathcal{O}_T)$ の対象
$$
K_i = Rf_{T, *}(E_{i, T} \otimes_{\mathcal{O}_{X_T}}^\mathbf{L} \mathcal{F})
$$
は perfect であり，その形成は基底変換と可換する；Derived Categories of Spaces, Lemma
\ref{spaces-perfect-lemma-base-change-tensor-perfect} を参照せよ．
したがって函数
$$
\chi_i : |T| \longrightarrow \mathbf{Z},\quad
\chi_i(t) =
\chi(X_t, E_{i, t} \otimes_{\mathcal{O}_{X_t}}^\mathbf{L} \mathcal{F}_t) =
\chi(K_i \otimes_{\mathcal{O}_T}^\mathbf{L} \kappa(t))
$$
は Derived Categories of Spaces, Lemma
\ref{spaces-perfect-lemma-chi-locally-constant} により局所定数である．
$P : I \to \mathbf{Z}$ を写像とする．数値的不変量が $P$ と一致する，
proper な台をもつ連接層の平坦族からなる部分スタック
$$
\Cohstack^P_{X/B} \subset \Cohstack_{X/B}
$$
を考える．より正確には，$\Cohstack_{X/B}$ の対象
$(T \to B, \mathcal{F})$ が $\Cohstack^P_{X/B}$ に属することと，
すべての $i \in I$ および $t \in T$ に対して $\chi_i(t) = P(i)$ となることは同値である．
\end{situation}

\begin{lemma}
\label{moduli-lemma-open-P}
Situation \ref{moduli-situation-numerical} において，スタック
$\Cohstack^P_{X/B}$ は代数的であり，
$$
\Cohstack^P_{X/B} \longrightarrow \Cohstack_{X/B}
$$
は平坦な閉埋め込みである．$I$ が有限であるか，または $B$ が局所 Noether 的ならば，
$\Cohstack^P_{X/B}$ は $\Cohstack_{X/B}$ の開かつ閉な部分スタックである．
\end{lemma}

\begin{proof}
$I$ が有限ならば，函数 $t \mapsto \chi_i(t)$ は局所定数なので，これは直ちに明らかである．
$I$ が無限ならば，
$$
I = \bigcup\nolimits_{I' \subset I\text{ finite}} I'
$$
と書き，$P' = P|_{I'}$ と記す．このとき
$$
\Cohstack^P_{X/B} = \bigcap\nolimits_{I' \subset I\text{ finite}}
\Cohstack^{P'}_{X/B}
$$
となる．したがって，$\Cohstack^P_{X/B}$ は常に代数スタックであり，射
$\Cohstack^P_{X/B} \subset \Cohstack_{X/B}$ は常に平坦な閉埋め込みであるが，
開部分スタックであるとは限らない（例の作成は読者に委ねる）．
しかし $B$ が局所 Noether 的ならば，Lemma \ref{moduli-lemma-coherent-qs-lfp} および
Morphisms of Stacks, Lemma
\ref{stacks-morphisms-lemma-locally-finite-type-locally-noetherian} により，
$\Cohstack_{X/B}$ も局所 Noether 的である．
したがって $U$ が局所 Noether スキームで，
$U \to \Cohstack_{X/B}$ が平滑かつ全射ならば，開かつ閉な部分スタック
$\Cohstack^{P'}_{X/B}$ の逆像どうしの共通部分は $U$ で開である
（局所 Noether 位相空間の連結成分は開だからである）．
これでこの場合の結果も従う．
\end{proof}

\begin{lemma}
\label{moduli-lemma-finite-list-perfect-objects}
$f : X \to B$ は本節の序文におけるものとする．
$E_1, \ldots, E_r \in D(\mathcal{O}_X)$ を perfect とする．
$I = \mathbf{Z}^{\oplus r}$ とおき，写像
$$
I \longrightarrow D(\mathcal{O}_X),\quad
(n_1, \ldots, n_r) \longmapsto
E_1^{\otimes n_1}
\otimes \ldots \otimes
E_r^{\otimes n_r}
$$
% P12-ERR-0196: Japanese supplies terminal sentence punctuation outside the frozen display.
を考える．$P : I \to \mathbf{Z}$ を写像とする．このとき，Situation
\ref{moduli-situation-numerical} で定義した
$\Cohstack^P_{X/B} \subset \Cohstack_{X/B}$ は開かつ閉な部分スタックである．
\end{lemma}

\begin{proof}
$B$ 上 \'etale 局所的に議論してよいので，$B$ はアフィンであると仮定してよい．
この場合には絶対 Noether 化による簡約を行える；読者には証明を飛ばすことを勧める．
実際，$B = \Spec(\Lambda)$ と書く．
各 $\Lambda_i$ が $\mathbf{Z}$ 上有限型となるフィルター付き余極限
$\Lambda = \colim \Lambda_i$ をとる．ある $i$ について，代数空間の射
$X_i \to \Spec(\Lambda_i)$ で分離的かつ有限表示であり，$\Lambda$ への基底変換が $X$ であるものをとれる．
Limits of Spaces, Lemmas
\ref{spaces-limits-lemma-descend-finite-presentation} および
\ref{spaces-limits-lemma-descend-separated-morphism} を参照せよ．
次に $i$ を大きくすれば，$D(\mathcal{O}_{X_i})$ の perfect 対象
$E_{1, i}, \ldots, E_{r, i}$ で，$X$ への導来引き戻しが
$E_1, \ldots, E_r$ と同型になるものが存在すると仮定できる；
Derived Categories of Spaces, Lemma
\ref{spaces-perfect-lemma-perfect-on-limit} を参照せよ．
明らかに，次の図式は Cartesian である：
$$
\xymatrix{
\Cohstack^P_{X/B} \ar[r] \ar[d] &
\Cohstack_{X/B} \ar[d] \\
\Cohstack^P_{X_i/\Spec(\Lambda_i)} \ar[r] &
\Cohstack_{X_i/\Spec(\Lambda_i)}
}
$$
したがって Lemma \ref{moduli-lemma-open-P} を適用すれば証明が完了する．
\end{proof}

\begin{example}[Hilbert 多項式を固定した連接層]
\label{moduli-example-hilbert-polynomial}
$f : X \to B$ は本節の序文におけるものとする．
$\mathcal{L}$ を可逆 $\mathcal{O}_X$-加群とする．
$P : \mathbf{Z} \to \mathbf{Z}$ を数値多項式とする．
このとき，数値的不変量が $P$ と一致する，proper な台をもつ連接層の平坦族からなる
開かつ閉な代数的部分スタック
$$
\Cohstack^P_{X/B} =
\Cohstack^{P, \mathcal{L}}_{X/B}
\subset \Cohstack_{X/B}
$$
を考えることができる．より正確には，$\Cohstack_{X/B}$ の対象
$(T \to B, \mathcal{F})$ が $\Cohstack^P_{X/B}$ に属することと，
$$
P(n) =
\chi(X_t, \mathcal{F}_t \otimes_{\mathcal{O}_{X_t}} \mathcal{L}_t^{\otimes n})
$$
がすべての $n \in \mathbf{Z}$ および $t \in T$ に対して成り立つことは同値である．
これはもちろん Situation \ref{moduli-situation-numerical} の特別な場合であり，そこでは
% P12-ERR-0197: The frozen malformed chain is retained rather than silently changing its mathematics.
$I = \mathbf{Z} \to D(\mathcal{O}_X)$ は
$n \mapsto \mathcal{L}^{\otimes n}$ によって与えられる．
Lemma \ref{moduli-lemma-finite-list-perfect-objects} から，これは開かつ閉な部分スタックである．
函数
$n \mapsto
\chi(X_t, \mathcal{F}_t \otimes_{\mathcal{O}_{X_t}} \mathcal{L}_t^{\otimes n})$
は常に数値多項式なので（Spaces over Fields, Lemma
\ref{spaces-over-fields-lemma-numerical-polynomial-from-euler}），次を得る：
$$
\Cohstack_{X/B} = \coprod\nolimits_{P\text{ numerical polynomial}}
\Cohstack^P_{X/B}
$$
すなわち，これは直和分解である．
\end{example}
% END EXPANDED MODULE ch108_01.tex
% BEGIN EXPANDED MODULE ch108_02.tex
\section{Quot の性質}
\label{moduli-section-quot}

\noindent
$f : X \to B$ を，分離的かつ有限表示な代数空間の射とする．
$\mathcal{F}$ を準連接 $\mathcal{O}_X$-加群とする．このとき
$\Quotfunctor_{\mathcal{F}/X/B}$ は代数空間である．
$\mathcal{F}$ が有限表示ならば，
$\Quotfunctor_{\mathcal{F}/X/B} \to B$ は局所有限表示である．
Quot, Proposition \ref{quot-proposition-quot} を参照せよ．

\begin{lemma}
\label{moduli-lemma-quot-diagonal-closed}
$\Quotfunctor_{\mathcal{F}/X/B} \to B$ の対角射は閉埋め込みである．
$\mathcal{F}$ が有限型ならば，その対角射は有限表示な閉埋め込みである．
\end{lemma}

\begin{proof}
スキーム $T/B$ と，$B$ 上の $\Quotfunctor_{\mathcal{F}/X/B}$ の
$T$-値点に対応する二つの商
$\mathcal{F}_T \to \mathcal{Q}_i$，$i = 1, 2$ が与えられたとする．
第一の商の核を $\mathcal{K}_1$ と記し，合成を
$u : \mathcal{K}_1 \to \mathcal{Q}_2$ とおく．
Flatness on Spaces, Lemma \ref{spaces-flat-lemma-F-zero-closed-proper} により，
$T$ のある閉部分空間が存在し，$T' \to T$ がそれを経由することと，
引き戻し $u_{T'}$ が零であることとは同値である．
これにより対角射が閉埋め込みであることが示された．
さらに $\mathcal{F}$ が有限型ならば，
$\mathcal{K}_1$ も有限型である
(Modules on Sites, Lemma
\ref{sites-modules-lemma-kernel-surjection-finite-onto-finite-presentation})．
同じ補題により，対角射が有限表示であることも分かる．
\end{proof}

\begin{lemma}
\label{moduli-lemma-quot-s-lfp}
射 $\Quotfunctor_{\mathcal{F}/X/B} \to B$ は分離的である．
$\mathcal{F}$ が有限表示ならば，この射はさらに局所有限表示である．
\end{lemma}

\begin{proof}
$\Quotfunctor_{\mathcal{F}/X/B} \to B$ が分離的であることを確認するには，
その対角射が閉埋め込みであることを示せばよい．これは Lemma
\ref{moduli-lemma-quot-diagonal-closed} により成り立つ．第二の主張は
Quot, Proposition \ref{quot-proposition-quot} の一部である．
\end{proof}

\begin{lemma}
\label{moduli-lemma-quot-existence-part}
$X \to B$ が proper かつ有限表示であり，$\mathcal{F}$ が有限型準連接加群であると仮定する．
このとき $\Quotfunctor_{\mathcal{F}/X/B} \to B$ は，付値判定法
(Morphisms of Spaces, Definition
\ref{spaces-morphisms-definition-valuative-criterion}) の存在部分を満たす．
\end{lemma}

\begin{proof}
基底変換を行えば，これは直ちに次の問題に帰着する：付値環 $R$ とその分数体 $K$，
$R$ 上 proper な代数空間 $X$，有限型準連接
$\mathcal{O}_X$-加群 $\mathcal{F}$，および連接な商
$\mathcal{F}_K \to \mathcal{Q}_K$ が与えられたとき，商
$\mathcal{F} \to \mathcal{Q}$ で，$\mathcal{Q}$ が有限表示
$\mathcal{O}_X$-加群であり，$R$ 上平坦かつ一般ファイバーが $\mathcal{Q}_K$ となるものの存在を示せ．
Flatness on Spaces, Theorem
\ref{spaces-flat-theorem-finite-type-flat} により，有限型準連接
$\mathcal{O}_X$-加群 $\mathcal{F}$ が $R$ 上平坦ならば有限表示であることに注意する．
まず $\mathcal{Q}$ の存在をアフィン局所的に解く．

\medskip\noindent
アフィン局所的には，次の問題に至る：
$R \to A$ を有限表示な環準同型とし，$M$ を有限 $A$-加群，
$\varphi : M_K \to N_K$ を $A_K$-加群の商とする．このとき
$$
L = \{x \in M \mid \varphi(x \otimes 1) = 0 \}
$$
を考えることができる．
% P12-ERR-0198: Japanese has no article corresponding to the frozen extraneous “The”.
$M \to M/L$ は $A$-加群の商であり，$R$-加群として torsion-free である．
したがって，これは $R$-加群として平坦である
(More on Algebra, Lemma
\ref{more-algebra-lemma-valuation-ring-torsion-free-flat})．
$M$ は有限 $A$-加群なので $L$ もそうであり，前掲の参照により
$L$ は有限表示 $A$-加群であると結論できる．
明らかに，$M/L$ は $(M/L)_K = N_K$ を満たすこのような商として一意である．

\medskip\noindent
前段落の構成における一意性により，これらの商は貼り合わされ，求める
$\mathcal{Q}$ を与える．もう少し詳しく述べよう．
全射かつ \'etale な射 $U \to X$ で，$U$ がアフィンスキームであるものを選ぶ．
上の構成を用いて，準連接で $R$ 上平坦，かつ一般ファイバー上で
$\mathcal{Q}_K|U$ を復元する商 $\mathcal{F}|_U \to \mathcal{Q}_U$ を構成する．
$X$ は分離的なので，$U \times_X U$ も $X$ 上 \'etale なアフィンスキームである．
このとき
$\mathcal{F}|_{U \times_X U} \to \text{pr}_1^*\mathcal{Q}_U$ と
$\mathcal{F}|_{U \times_X U} \to \text{pr}_2^*\mathcal{Q}_U$ は，
% P12-ERR-0199: Japanese renders the intended uniqueness despite the frozen typo “uniquess”.
構成の一意性により商として一致する．したがって，求める全射
$\mathcal{F} \to \mathcal{Q}$ へ $\mathcal{F}|_U \to \mathcal{Q}_U$ を降下できる
(Properties of Spaces, Proposition
\ref{spaces-properties-proposition-quasi-coherent})．
\end{proof}

\begin{lemma}
\label{moduli-lemma-quot-functorial}
$B$ を代数空間とする．$\pi : X \to Y$ を，$B$ 上分離的かつ有限表示な
代数空間の間のアフィンかつ準有限な射とする．
$\mathcal{F}$ を準連接 $\mathcal{O}_X$-加群とする．このとき $\pi_*$ は射
$\Quotfunctor_{\mathcal{F}/X/B} \to \Quotfunctor_{\pi_*\mathcal{F}/Y/B}$ を誘導する．
\end{lemma}

\begin{proof}
$\mathcal{G} = \pi_*\mathcal{F}$ とおく．$\pi$ はアフィンなので，$B$ 上の任意のスキーム
$T$ に対して，Cohomology of Spaces, Lemma
\ref{spaces-cohomology-lemma-affine-base-change} により
$\mathcal{G}_T = \pi_{T, *}\mathcal{F}_T$ である．
さらに $\pi_T$ はアフィンなので，$\pi_{T, *}$ は完全であり，商を商へ移す．
準連接な商 $\mathcal{F}_T \to \mathcal{Q}$ が
% P12-ERR-0200: The frozen Quot functor missing the sheaf parameter is retained literally.
$\Quotfunctor_{X/B}$ の点を定めることと，$\mathcal{Q}$ が $T$ 上の
$\Cohstack_{X/B}$ の対象を定めることとは同値である
（$\mathcal{G}$ と $Y$ についても同様である）．
Lemma \ref{moduli-lemma-coherent-functorial} で $\pi_*$ が射
$\Cohstack_{X/B} \to \Cohstack_{Y/B}$ を誘導することを見たので，
$\mathcal{F}_T \to \mathcal{Q}$ が
$\Quotfunctor_{\mathcal{F}/X/B}(T)$ に属するならば，
$\mathcal{G}_T \to \pi_{T, *}\mathcal{Q}$ は
$\Quotfunctor_{\mathcal{G}/Y/B}(T)$ に属する．
\end{proof}

\begin{lemma}
\label{moduli-lemma-quot-open}
$B$ を代数空間とする．$\pi : X \to Y$ を，$B$ 上分離的かつ有限表示な
代数空間の間のアフィンな開埋め込みとする．
$\mathcal{F}$ を準連接 $\mathcal{O}_X$-加群とする．このとき Lemma
\ref{moduli-lemma-quot-functorial} の射
$\Quotfunctor_{\mathcal{F}/X/B} \to \Quotfunctor_{\pi_*\mathcal{F}/Y/B}$ は開埋め込みである．
\end{lemma}

\begin{proof}
省略する．ヒント：$(\pi_*\mathcal{F})_T \to \mathcal{Q}$ が
$\Quotfunctor_{\pi_*\mathcal{F}/Y/B}(T)$ の元であり，
$t \in T$ に対して $\text{Supp}(\mathcal{Q}_t) \subset |X_t|$ ならば，
$t$ のある近傍に属する $t' \in T$ についても同じことが成り立つ．
\end{proof}

\begin{lemma}
\label{moduli-lemma-quot-better-open}
$B$ を代数空間とする．$j : X \to Y$ を，$B$ 上分離的かつ有限表示な
代数空間の開埋め込みとする．$\mathcal{G}$ を準連接 $\mathcal{O}_Y$-加群とし，
$\mathcal{F} = j^*\mathcal{G}$ とおく．このとき，$B$ 上の代数空間の開埋め込み
$$
\Quotfunctor_{\mathcal{F}/X/B}
\longrightarrow
\Quotfunctor_{\mathcal{G}/Y/B}
$$
が存在する．
\end{lemma}

\begin{proof}
$\mathcal{F}_T \to \mathcal{Q}$ が
$\Quotfunctor_{\mathcal{F}/X/B}(T)$ の元ならば，合成
$\mathcal{G}_T \to j_{T, *}\mathcal{F}_T \to j_{T, *}\mathcal{Q}$ を考えることができる．
茎を調べると，これは全射であると分かる．
Lemma \ref{moduli-lemma-coherent-functorial} により，$j_{T, *}\mathcal{Q}$ は有限表示で，
% P12-ERR-0201: Both frozen base-B assertions are retained rather than silently changed to T.
$B$ 上平坦，かつ台が $B$ 上 proper である．したがって
$\Quotfunctor_{\mathcal{G}/Y/B}$ の $T$-値点を得る．
これが補題の射を定める．この射が開埋め込みであることの証明は省略する．ヒント：
$\mathcal{G}_T \to \mathcal{Q}$ が
$\Quotfunctor_{\mathcal{G}/Y/B}(T)$ の元であり，
$t \in T$ に対して $\text{Supp}(\mathcal{Q}_t) \subset |X_t|$ ならば，
$t$ のある近傍に属する $t' \in T$ についても同じことが成り立つ．
\end{proof}


\begin{lemma}
\label{moduli-lemma-quot-closed}
$B$ を代数空間とする．$\pi : X \to Y$ を，$B$ 上分離的かつ有限表示な
代数空間の閉埋め込みとする．$\mathcal{F}$ を準連接 $\mathcal{O}_X$-加群とする．
このとき Lemma \ref{moduli-lemma-quot-functorial} の射
$\Quotfunctor_{\mathcal{F}/X/B} \to \Quotfunctor_{\pi_*\mathcal{F}/Y/B}$ は同型である．
\end{lemma}

\begin{proof}
$B$ 上の任意のスキーム $T$ に対して，射 $\pi_T : X_T \to Y_T$ は閉埋め込みである．
このとき $\pi_{T, *}$ は，$\QCoh(\mathcal{O}_{X_T})$ と，
$X_T$ のイデアル層によって零化される準連接加群を対象とする
$\QCoh(\mathcal{O}_{Y_T})$ の充満部分圏との間の圏同値である；
Morphisms of Spaces, Lemma \ref{spaces-morphisms-lemma-i-star-equivalence} を参照せよ．
% P12-ERR-0202: Japanese renders the intended quotient despite the frozen typo “qotient”.
$(\pi_*\mathcal{F})_T$ の商はこのイデアルによって零化されるので，任意の $T$ に対する写像
$\Quotfunctor_{\mathcal{F}/X/B}(T) \to \Quotfunctor_{\pi_*\mathcal{F}/Y/B}(T)$
の全単射性が得られ，主張が従う．
\end{proof}

\begin{lemma}
\label{moduli-lemma-quot-quotient}
$X \to B$ は本節の序文におけるものとする．準連接 $\mathcal{O}_X$-加群の全射
$\mathcal{F} \to \mathcal{G}$ をとる．このとき標準的な閉埋め込み
$\Quotfunctor_{\mathcal{G}/X/B} \to \Quotfunctor_{\mathcal{F}/X/B}$ が存在する．
\end{lemma}

\begin{proof}
$\mathcal{K} = \Ker(\mathcal{F} \to \mathcal{G})$ とおく．
引き戻しの右完全性により，$B$ 上の任意のスキーム $T$ に対して
$\mathcal{K}_T \to \mathcal{F}_T \to \mathcal{G}_T \to 0$ は
% P12-ERR-0203: Japanese renders the intended exact sequence despite the frozen typo “sequecnce”.
完全列である．特に $\mathcal{G}_T$ の商は $\mathcal{F}_T$ の商を定めるので，関手の変換
$\Quotfunctor_{\mathcal{G}/X/B} \to \Quotfunctor_{\mathcal{F}/X/B}$ を得る．
この変換は Flatness on Spaces, Lemma
\ref{spaces-flat-lemma-F-zero-closed-proper} により閉埋め込みである．
% P12-ERR-0204: Japanese repairs only the frozen English sentence grammar.
実際，$\Quotfunctor_{\mathcal{F}/X/B}(T)$ の元
$\mathcal{F}_T \to \mathcal{Q}$ が与えられたとき，$T'/T$ への引き戻しが
この変換の像に属することと，$\mathcal{K}_{T'} \to \mathcal{Q}_{T'}$ が零であることとは同値である．
\end{proof}

\begin{remark}[数値的不変量]
\label{moduli-remark-quot-numerical}
$f : X \to B$ および $\mathcal{F}$ は本節の序文におけるものとする．
$I$ を集合とし，各 $i \in I$ に対して $E_i \in D(\mathcal{O}_X)$ を perfect とする．
$P : I \to \mathbf{Z}$ を函数とする．次の射があることを想起せよ：
$$
\Quotfunctor_{\mathcal{F}/X/B} \longrightarrow \Cohstack_{X/B}
$$
これは $\Quotfunctor_{\mathcal{F}/X/B}(T)$ の元
$\mathcal{F}_T \to \mathcal{Q}$ を，$T$ 上の $\Cohstack_{X/B}$ の対象
$\mathcal{Q}$ へ送る；Quot, Proposition \ref{quot-proposition-quot} の証明を参照せよ．
したがって，ファイバー積の図式
$$
\xymatrix{
\Quotfunctor^P_{\mathcal{F}/X/B} \ar[r] \ar[d] &
\Cohstack^P_{X/B} \ar[d] \\
\Quotfunctor_{\mathcal{F}/X/B} \ar[r] &
\Cohstack_{X/B}
}
$$
を作れる．これは左上隅の代数空間を定める図式である．左の垂直射は平坦な閉埋め込みであり，
例えば $I$ が有限であるか，$B$ が局所 Noether 的であるか，または
$I = \mathbf{Z}$ かつ，ある可逆 $\mathcal{O}_X$-加群 $\mathcal{L}$ に対して
$E_i = \mathcal{L}^{\otimes i}$ である場合には，開かつ閉な埋め込みである
（最後の場合には記法 $\Quotfunctor^{P, \mathcal{L}}_{\mathcal{F}/X/B}$ を用いることがある）．
Situation \ref{moduli-situation-numerical}，Lemmas \ref{moduli-lemma-open-P} および
\ref{moduli-lemma-finite-list-perfect-objects}，ならびに Example
\ref{moduli-example-hilbert-polynomial} を参照せよ．
\end{remark}

\begin{lemma}
\label{moduli-lemma-quot-tensor-invertible}
$f : X \to B$ および $\mathcal{F}$ は本節の序文におけるものとする．
$\mathcal{L}$ を可逆 $\mathcal{O}_X$-加群とする．このとき $\mathcal{L}$ とのテンソル積は同型
$$
\Quotfunctor_{\mathcal{F}/X/B}
\longrightarrow
\Quotfunctor_{\mathcal{F} \otimes_{\mathcal{O}_X} \mathcal{L}/X/B}
$$
を定める．
% P12-ERR-0205: Japanese repairs only the frozen English “Given ..., then ... this map” grammar.
数値多項式 $P(t)$ が与えられたとき，$P'(t) = P(t + 1)$ とおけば，この写像は
開かつ閉な部分スタックの同型
$\Quotfunctor^P_{\mathcal{F}/X/B}
\longrightarrow
\Quotfunctor^{P'}_{\mathcal{F} \otimes_{\mathcal{O}_X} \mathcal{L}/X/B}$
を誘導する．
\end{lemma}

\begin{proof}
$\mathcal{G} = \mathcal{F} \otimes_{\mathcal{O}_X} \mathcal{L}$ とおく．
$\mathcal{G}_T = \mathcal{F}_T \otimes_{\mathcal{O}_{X_T}} \mathcal{L}_T$ であることに注意する．
$\mathcal{F}_T \to \mathcal{Q}$ が
$\Quotfunctor_{\mathcal{F}/X/B}(T)$ の元ならば，これを
$\Quotfunctor_{\mathcal{F} \otimes_{\mathcal{O}_X} \mathcal{L}/X/B}(T)$ の元
$\mathcal{G}_T \to \mathcal{Q} \otimes_{\mathcal{O}_{X_T}} \mathcal{L}_T$ へ送る．
これは引き戻しと可換するので，求める関手の変換を定める．
明らかな逆変換が存在するから，同型である．最後の主張の証明は省略する．
\end{proof}

\begin{lemma}
\label{moduli-lemma-quot-power-invertible}
$f : X \to B$ および $\mathcal{F}$ は本節の序文におけるものとする．
$\mathcal{L}$ を可逆 $\mathcal{O}_X$-加群とする．このとき
$$
\Quotfunctor^{P, \mathcal{L}}_{\mathcal{F}/X/B} =
\Quotfunctor^{P', \mathcal{L}^{\otimes n}}_{\mathcal{F}/X/B}
$$
である．ここで $P'(t) = P(nt)$ とする．
\end{lemma}

\begin{proof}
すべての定義を展開すれば直ちに従う．
\end{proof}
% END EXPANDED MODULE ch108_02.tex
% BEGIN EXPANDED MODULE ch108_03.tex
\section{Quot に対する有界性}
\label{moduli-section-quot-bounded}

\noindent
古典的な場合とは異なり，Quot 関手が代数空間であることは既に分かっているが，
それが有限型代数空間によって表現される場合があるかどうかは，まだ分かっていない．

\begin{lemma}
\label{moduli-lemma-quot-Pn}
$n \geq 0$，$r \geq 1$，$P \in \mathbf{Q}[t]$ とする．
Hilbert 多項式が $P$ である $\mathcal{O}_{\mathbf{P}^n_\mathbf{Z}}^{\oplus r}$ の商を
パラメータ付ける代数空間
$$
X = \Quotfunctor^P_{\mathcal{O}^{\oplus r}_{\mathbf{P}^n_\mathbf{Z}}/
\mathbf{P}^n_\mathbf{Z}/\mathbf{Z}}
$$
は $\Spec(\mathbf{Z})$ 上 proper である．
\end{lemma}

\begin{proof}
$X \to \Spec(\mathbf{Z})$ が分離的かつ局所有限表示であることは既に分かっている
(Lemma \ref{moduli-lemma-quot-s-lfp})．また，Lemma
\ref{moduli-lemma-quot-existence-part} により，$X \to \Spec(\mathbf{Z})$ は
付値判定法の存在部分を満たすことも分かっている．
proper 性の付値判定法により，Quot 空間が準コンパクトであることを示せば十分である；
Morphisms of Spaces, Lemma
\ref{spaces-morphisms-lemma-characterize-proper} を参照せよ．
したがって，準コンパクトスキーム $T$ と全射 $T \to X$ を見つければ十分である．
Varieties, Lemma \ref{varieties-lemma-bound-quotients-free} で得られる整数を $m$ とする．
次のようにおく：
$$
N = r{m + n \choose n} - P(m)
$$
以下，$\mathbf{P}^n$ は
$\mathbf{P}^n_\mathbf{Z} = \text{Proj}(\mathbf{Z}[T_0, \ldots, T_n])$ を表すものとし，
添字のない積は $\Spec(\mathbf{Z})$ 上の積を意味するものとする．
証明の考え方は，アフィンスキーム $T$ 上で「普遍的」な写像
$$
\Psi :
\mathcal{O}_{T \times \mathbf{P}^n}(-m)^{\oplus N}
\longrightarrow
\mathcal{O}_{T \times \mathbf{P}^n}^{\oplus r}
$$
を構成し，$X$ の各点が $T$ のある点におけるこの写像の余核に対応することを示すことである．

\medskip\noindent
$T$ と $\Psi$ の定義．$T = \Spec(A)$ とし，ここで
$$
A = \mathbf{Z}[a_{i, j, E}]
$$
とする．添字は $i \in \{1, \ldots, r\}$，$j \in \{1, \ldots, N\}$，
および全次数
% P12-ERR-0206: The frozen missing comma in the summation range is retained literally.
$|E| = \sum_{k = 0, \ldots n} e_k = m$ をもつ多重指数
$E = (e_0, \ldots, e_n)$ を走る．
次に，$\Psi$ を，その $(i, j)$ 行列成分が写像
$$
\sum\nolimits_{E = (e_0, \ldots, e_n)}
a_{i, j, E} T_0^{e_0} \ldots T_n^{e_n} :
\mathcal{O}_{T \times \mathbf{P}^n}(-m)
\longrightarrow
\mathcal{O}_{T \times \mathbf{P}^n}
$$
であるものとして定める．ここで和は上記の $E$ を走る（もちろん $i$ と $j$ は固定している）．

\medskip\noindent
$T \times \mathbf{P}^n$ 上の商 $\mathcal{Q} = \Coker(\Psi)$ を考える．
More on Morphisms, Lemma
\ref{more-morphisms-lemma-generic-flatness-stratification} により，ある $t \geq 0$ と閉部分スキーム
$$
T = T_0 \supset T_1 \supset \ldots \supset T_t = \emptyset
$$
が存在し，$\mathcal{Q}$ の $(T_p \setminus T_{p + 1}) \times \mathbf{P}^n$ への引き戻し
$\mathcal{Q}_p$ は $T_p \setminus T_{p + 1}$ 上平坦である．
$\mathcal{Q} = \Coker(\Psi)$ を定める完全列を引き戻すことにより，完全列
$$
\mathcal{O}_{(T_p \setminus T_{p + 1}) \times \mathbf{P}^n}(-m)^{\oplus N}
\to
\mathcal{O}_{(T_p \setminus T_{p + 1}) \times \mathbf{P}^n}^{\oplus r}
\to
\mathcal{Q}_p
\to
0
$$
を得ることに注意する．したがって射
$$
% P12-ERR-0207: The frozen coproduct lacks its index and range; it is retained literally.
\coprod (T_p \setminus T_{p + 1})
\longrightarrow
% P12-ERR-0208: Both frozen projective-space denominators omit the exponent n.
\Quotfunctor_{\mathcal{O}^{\oplus r}/\mathbf{P}/\mathbf{Z}}
\supset
\Quotfunctor^P_{\mathcal{O}^{\oplus r}/\mathbf{P}/\mathbf{Z}} = X
$$
を得る．左辺は Noether スキームであり，右辺の包含は開なので，
$X$ の任意の点がこの射の像に含まれることを示せば十分である．

\medskip\noindent
$k$ を体とし，$x \in X(k)$ とする．このとき $x$ は，Hilbert 多項式が $P$ である
連接 $\mathcal{O}_{\mathbf{P}^n_k}$-加群 $\mathcal{F}$ への全射
$\mathcal{O}_{\mathbf{P}^n_k}^{\oplus r} \to \mathcal{F}$ に対応する．
短完全列
$$
0 \to \mathcal{K} \to
\mathcal{O}_{\mathbf{P}^n_k}^{\oplus r} \to
\mathcal{F} \to 0
$$
を考える．Varieties, Lemma \ref{varieties-lemma-bound-quotients-free} と
$m$ の選び方により，$\mathcal{K}$ は $m$-regular である．
Varieties, Lemma \ref{varieties-lemma-m-regular-globally-generated} により
$\mathcal{K}(m)$ は大域生成される．Varieties, Lemma
\ref{varieties-lemma-m-regular-up} と $m$-regularity の定義により，
$i > 0$ に対して $H^i(\mathbf{P}^n_k, \mathcal{K}(m)) = 0$ である．
ゆえに
$$
\dim_k H^0(\mathbf{P}^n_k, \mathcal{K}(m)) =
\chi(\mathcal{K}(m)) =
\chi(\mathcal{O}_{\mathbf{P}^n_k}(m)^{\oplus r}) -
\chi(\mathcal{F}(m)) = N
$$
となることが，$N$ の選び方から分かる．これにより全射
$$
\mathcal{O}_{\mathbf{P}^n_k}^{\oplus N}
\longrightarrow
\mathcal{K}(m)
$$
を得る．再び twist を戻し，上の短完全列を用いると，$\mathcal{F}$ は写像
$$
\Psi_x :
\mathcal{O}_{\mathbf{P}^n_k}(-m)^{\oplus N}
\to
\mathcal{O}_{\mathbf{P}^n_k}^{\oplus r}
$$
の余核であることが分かる．対応する射 $t = \Spec(\tau) : \Spec(k) \to T$ による
$\Psi$ の基底変換が $\Psi_x$ となるような環準同型 $\tau : A \to k$ が一意に存在する．
実際，$\Psi_x$ を定める
% P12-ERR-0209: The frozen N-by-r matrix dimension is retained literally.
$N \times r$ 行列の成分は，$T_0, \ldots, T_n$ における次数 $m$ の斉次多項式
$\sum \lambda_{i, j, E} T_0^{e_0} \ldots T_n^{e_n}$ であり，係数は
$\lambda_{i, j, E} \in k$ であるから，
$\tau(a_{i, j, E}) = \lambda_{i, j, E}$ とおけばよい．
すると，ある $p$ に対して $t \in T_p \setminus T_{p + 1}$ であり，
上の射による $t$ の像は求めるとおり $x$ である．
\end{proof}

\begin{lemma}
\label{moduli-lemma-quot-Pn-over-base}
$B$ を代数空間とする．$X = B \times \mathbf{P}^n_\mathbf{Z}$ とおく．
$\mathcal{L}$ を $\mathcal{O}_{\mathbf{P}^n}(1)$ の $X$ への引き戻しとする．
$\mathcal{F}$ を有限表示 $\mathcal{O}_X$-加群とする．
$\mathcal{F}$ の商で，$\mathcal{L}$ に関する Hilbert 多項式が $P$ であるものを
パラメータ付ける代数空間 $\Quotfunctor^P_{\mathcal{F}/X/B}$ は $B$ 上 proper である．
\end{lemma}

\begin{proof}
この問題は $B$ 上 \'etale 局所的である；Morphisms of Spaces, Lemma
\ref{spaces-morphisms-lemma-proper-local} を参照せよ．
したがって，$B$ はアフィンスキームであると仮定してよい．
この場合，$\mathcal{L}$ は $X$ 上豊富な可逆加群である
(Constructions, Lemma \ref{constructions-lemma-ample-on-proj} および
Properties, Definition \ref{properties-definition-ample} の豊富な可逆加群の定義による)．
したがって，ある $r' \geq 0$ と $r \geq 0$，および全射
$$
\mathcal{O}_X^{\oplus r} \longrightarrow
\mathcal{F} \otimes_{\mathcal{O}_X} \mathcal{L}^{\otimes r'}
$$
が存在する；Properties, Proposition
\ref{properties-proposition-characterize-ample} を参照せよ．
Lemma \ref{moduli-lemma-quot-tensor-invertible} により，$\mathcal{F}$ を
$\mathcal{F} \otimes_{\mathcal{O}_X} \mathcal{L}^{\otimes r'}$ に，
$P(t)$ を $P(t + r')$ に置き換えてよい．
Lemma \ref{moduli-lemma-quot-quotient} により閉埋め込み
$$
\Quotfunctor^P_{\mathcal{F}/X/B}
\longrightarrow
\Quotfunctor^P_{\mathcal{O}_X^{\oplus r}/X/B}
$$
を得る．Lemma \ref{moduli-lemma-quot-Pn} で
$\Quotfunctor^P_{\mathcal{O}_X^{\oplus r}/X/B} \to B$ が proper であることを示したので，結論を得る．
\end{proof}

\begin{lemma}
\label{moduli-lemma-quot-proper-over-base}
$f : X \to B$ を，代数空間の proper かつ有限表示な射とする．
$\mathcal{F}$ を有限表示 $\mathcal{O}_X$-加群とする．
$\mathcal{L}$ を $X/B$ 上豊富な可逆 $\mathcal{O}_X$-加群とする；
Divisors on Spaces, Definition
\ref{spaces-divisors-definition-relatively-ample} を参照せよ．
$\mathcal{F}$ の商で，$\mathcal{L}$ に関する Hilbert 多項式が $P$ であるものを
パラメータ付ける代数空間 $\Quotfunctor^P_{\mathcal{F}/X/B}$ は $B$ 上 proper である．
\end{lemma}

\begin{proof}
この問題は $B$ 上 \'etale 局所的である；Morphisms of Spaces, Lemma
\ref{spaces-morphisms-lemma-proper-local} を参照せよ．
したがって，$B$ はアフィンスキームであると仮定してよい．
このとき，ある閉埋め込み $i : X \to \mathbf{P}^n_B$ と $d \geq 1$ が存在して，
$i^*\mathcal{O}_{\mathbf{P}^n_B}(1) \cong \mathcal{L}^{\otimes d}$ となる．
Morphisms, Lemma \ref{morphisms-lemma-quasi-projective-finite-type-over-S} を参照せよ．
$\mathcal{L}$ を $\mathcal{L}^{\otimes d}$ に，数値多項式 $P(t)$ を $P(dt)$ に置き換えても
$\Quotfunctor^P_{\mathcal{F}/X/B}$ は変わらない（細部は省略する）．
したがって $\mathcal{L} = i^*\mathcal{O}_{\mathbf{P}^n_B}(1)$ と仮定してよい．
Lemma \ref{moduli-lemma-quot-closed} の同型
$\Quotfunctor_{\mathcal{F}/X/B} \to
\Quotfunctor_{i_*\mathcal{F}/\mathbf{P}^n_B/B}$ は同型
$\Quotfunctor^P_{\mathcal{F}/X/B} \cong
\Quotfunctor^P_{i_*\mathcal{F}/\mathbf{P}^n_B/B}$ を誘導する．
Lemma \ref{moduli-lemma-quot-Pn-over-base} により
$\Quotfunctor^P_{i_*\mathcal{F}/\mathbf{P}^n_B/B}$ は $B$ 上 proper なので，結論を得る．
\end{proof}

\begin{lemma}
\label{moduli-lemma-quot-qc-over-base}
$f : X \to B$ を，代数空間の分離的かつ有限表示な射とする．
$\mathcal{F}$ を有限表示 $\mathcal{O}_X$-加群とする．
$\mathcal{L}$ を $X/B$ 上豊富な可逆 $\mathcal{O}_X$-加群とする；
Divisors on Spaces, Definition
\ref{spaces-divisors-definition-relatively-ample} を参照せよ．
$\mathcal{F}$ の商で，$\mathcal{L}$ に関する Hilbert 多項式が $P$ であるものを
パラメータ付ける代数空間 $\Quotfunctor^P_{\mathcal{F}/X/B}$ は，
$B$ 上分離的かつ有限表示である．
\end{lemma}

\begin{proof}
$\Quotfunctor_{\mathcal{F}/X/B} \to B$ が分離的かつ局所有限表示であることは既に見た；
Lemma \ref{moduli-lemma-quot-s-lfp} を参照せよ．したがって，Remark
\ref{moduli-remark-quot-numerical} の開部分空間
$\Quotfunctor^P_{\mathcal{F}/X/B}$ が $B$ 上準コンパクトであることを示せば十分である．

\medskip\noindent
この問題は $B$ 上 \'etale 局所的である
(Morphisms of Spaces, Lemma \ref{spaces-morphisms-lemma-quasi-compact-local})．
したがって $B$ はアフィンであると仮定してよい．

\medskip\noindent
$B = \Spec(\Lambda)$ と仮定する．この環を有限型 $\mathbf{Z}$-部分代数の余極限
$\Lambda = \colim \Lambda_i$ として書く．このとき，ある $i$ と，
$B_i = \Spec(\Lambda_i)$ 上で補題と同じ条件を満たす系
$X_i, \mathcal{F}_i, \mathcal{L}_i$ で，$B$ への基底変換が
$X, \mathcal{F}, \mathcal{L}$ を与えるものをとれる．これは Limits of Spaces の次の補題から従う：
\ref{spaces-limits-lemma-descend-finite-presentation}（$X_i$ を得るため），
\ref{spaces-limits-lemma-descend-modules-finite-presentation}（$\mathcal{F}_i$ を得るため），
\ref{spaces-limits-lemma-descend-invertible-modules}（$\mathcal{L}_i$ を得るため），および
\ref{spaces-limits-lemma-descend-separated}（$X_i$ を分離的にするため）．
次が成り立つので，
$$
\Quotfunctor_{\mathcal{F}/X/B} = B \times_{B_i}
\Quotfunctor_{\mathcal{F}_i/X_i/B_i}
$$
$\Quotfunctor^P_{\mathcal{F}/X/B}$ についても同様に，次の段落で扱う場合に帰着する．

\medskip\noindent
$B$ はアフィンかつ Noether 的であると仮定する．Lemma
\ref{moduli-lemma-quot-power-invertible} により，$\mathcal{L}$ を正の冪に置き換えてよい．
したがって，ある immersion $i : X \to \mathbf{P}^n_B$ が存在して
$i^*\mathcal{O}_{\mathbf{P}^n}(1) = \mathcal{L}$ となると仮定してよい．
Morphisms, Lemma \ref{morphisms-lemma-quasi-compact-immersion} により，
$i$ が開埋め込み $j : X \to X'$ を経由するような閉部分スキーム
$X' \subset \mathbf{P}^n_B$ が存在する．Properties, Lemma
\ref{properties-lemma-lift-finite-presentation} により，有限表示
$\mathcal{O}_{X'}$-加群 $\mathcal{G}$ で $j^*\mathcal{G} = \mathcal{F}$ となるものが存在する．
したがって Lemma \ref{moduli-lemma-quot-better-open} により開埋め込み
$$
\Quotfunctor_{\mathcal{F}/X/B}
\longrightarrow
\Quotfunctor_{\mathcal{G}/X'/B}
$$
を得る．明らかに，この開埋め込みは $\Quotfunctor^P_{\mathcal{F}/X/B}$ を
$\Quotfunctor^P_{\mathcal{G}/X'/B}$ の中へ送る．ところで
Lemma \ref{moduli-lemma-quot-proper-over-base} により
$\Quotfunctor^P_{\mathcal{G}/X'/B}$ は $B$ 上 proper である．
したがってこれは Noether 的であり，Noether 的代数空間の任意の開部分は準コンパクトなので，証明が完了する．
\end{proof}
% END EXPANDED MODULE ch108_03.tex
% BEGIN EXPANDED MODULE ch108_04.tex
\section{Hilbert 関手の性質}
\label{moduli-section-hilb}

\noindent
$f : X \to B$ を，分離的かつ有限表示な代数空間の射とする．このとき
$\Hilbfunctor_{X/B}$ は $B$ 上局所有限表示な代数空間である．
Quot, Proposition \ref{quot-proposition-hilb} を参照せよ．

\begin{lemma}
\label{moduli-lemma-hilb-diagonal-closed}
$\Hilbfunctor_{X/B} \to B$ の対角射は有限表示な閉埋め込みである．
\end{lemma}

\begin{proof}
Quot, Lemma \ref{quot-lemma-hilb-is-quot} で
$\Hilbfunctor_{X/B} = \Quotfunctor_{\mathcal{O}_X/X/B}$ であることを見た．
したがって Lemma \ref{moduli-lemma-quot-diagonal-closed} から従う．
\end{proof}

\begin{lemma}
\label{moduli-lemma-hilb-s-lfp}
射 $\Hilbfunctor_{X/B} \to B$ は分離的かつ局所有限表示である．
\end{lemma}

\begin{proof}
$\Hilbfunctor_{X/B} \to B$ が分離的であることを確認するには，
その対角射が閉埋め込みであることを示せばよい．これは Lemma
\ref{moduli-lemma-hilb-diagonal-closed} により成り立つ．第二の主張は Quot, Proposition
\ref{quot-proposition-hilb} の一部である．
\end{proof}

\begin{lemma}
\label{moduli-lemma-hilb-existence-part}
$X \to B$ が proper かつ有限表示であると仮定する．このとき
$\Hilbfunctor_{X/B} \to B$ は，付値判定法
(Morphisms of Spaces, Definition
\ref{spaces-morphisms-definition-valuative-criterion}) の存在部分を満たす．
\end{lemma}

\begin{proof}
Quot, Lemma \ref{quot-lemma-hilb-is-quot} で
$\Hilbfunctor_{X/B} = \Quotfunctor_{\mathcal{O}_X/X/B}$ であることを見た．
したがって Lemma \ref{moduli-lemma-quot-existence-part} から従う．
\end{proof}

\begin{lemma}
\label{moduli-lemma-hilb-open}
$B$ を代数空間とする．$\pi : X \to Y$ を，$B$ 上分離的かつ有限表示な
代数空間の開埋め込みとする．このとき $\pi$ は開埋め込み
$\Hilbfunctor_{X/B} \to \Hilbfunctor_{Y/B}$ を誘導する．
\end{lemma}

\begin{proof}
省略する．ヒント：$T$ 上 proper な閉部分スキーム $Z \subset X_T$ が与えられたとき，
$Z$ は $Y_T$ においても閉である．したがって変換
$\Hilbfunctor_{X/B} \to \Hilbfunctor_{Y/B}$ を得る．
$Z \subset Y_T$ が $\Hilbfunctor_{Y/B}(T)$ の元であり，
$t \in T$ に対して $|Z_t| \subset |X_t|$ ならば，
$t$ のある近傍に属する $t' \in T$ についても同じことが成り立つ．
\end{proof}

\begin{lemma}
\label{moduli-lemma-hilb-closed}
$B$ を代数空間とする．$\pi : X \to Y$ を，$B$ 上分離的かつ有限表示な
代数空間の閉埋め込みとする．このとき $\pi$ は閉埋め込み
$\Hilbfunctor_{X/B} \to \Hilbfunctor_{Y/B}$ を誘導する．
\end{lemma}

\begin{proof}
$\pi$ は閉埋め込みなので，閉部分スキーム $Z \subset X_T$ が与えられたとき，
% P12-ERR-0210: The frozen tautological target X_T is retained rather than silently changed to Y_T.
$Z$ を $X_T$ の閉部分スキームとみなせることは直ちに分かる．したがって変換
$\Hilbfunctor_{X/B} \to \Hilbfunctor_{Y/B}$ を得る．
この変換が単射であることは直ちに分かる．これが閉埋め込みであることを示すには，写像
$\mathcal{O}_Y \to \mathcal{O}_X$ に Lemma \ref{moduli-lemma-quot-quotient} を用い，
Quot, Lemma \ref{quot-lemma-hilb-is-quot} の同一視
$\Hilbfunctor_{X/B} = \Quotfunctor_{\mathcal{O}_X/X/B}$，
$\Hilbfunctor_{Y/B} = \Quotfunctor_{\mathcal{O}_Y/Y/B}$ を用いればよい．
\end{proof}

\begin{remark}[数値的不変量]
\label{moduli-remark-hilb-numerical}
$f : X \to B$ は本節の序文におけるものとする．
$I$ を集合とし，各 $i \in I$ に対して $E_i \in D(\mathcal{O}_X)$ を perfect とする．
$P : I \to \mathbf{Z}$ を函数とする．Quot, Lemma
\ref{quot-lemma-hilb-is-quot} により
$\Hilbfunctor_{X/B} = \Quotfunctor_{\mathcal{O}_X/X/B}$ であることを想起せよ．
したがって，
$$
\Hilbfunctor^P_{X/B} = \Quotfunctor^P_{\mathcal{O}_X/X/B}
$$
と定義できる．ここで $\Quotfunctor^P_{\mathcal{O}_X/X/B}$ は Remark
\ref{moduli-remark-quot-numerical} におけるものである．射
$$
\Hilbfunctor^P_{X/B} \longrightarrow \Hilbfunctor_{X/B}
$$
は平坦な閉埋め込みであり，例えば $I$ が有限であるか，$B$ が局所 Noether 的であるか，または
$I = \mathbf{Z}$ かつ，ある可逆 $\mathcal{O}_X$-加群 $\mathcal{L}$ に対して
$E_i = \mathcal{L}^{\otimes i}$ である場合には，開かつ閉な埋め込みである．
最後の場合には，記法 $\Hilbfunctor^{P, \mathcal{L}}_{X/B}$ を用いることがある．
\end{remark}

\begin{lemma}
\label{moduli-lemma-hilb-proper-over-base}
$f : X \to B$ を，代数空間の proper かつ有限表示な射とする．
$\mathcal{L}$ を $X/B$ 上豊富な可逆 $\mathcal{O}_X$-加群とする；
Divisors on Spaces, Definition
\ref{spaces-divisors-definition-relatively-ample} を参照せよ．
$\mathcal{L}$ に関する Hilbert 多項式が $P$ である閉部分スキームを
パラメータ付ける代数空間 $\Hilbfunctor^P_{X/B}$ は $B$ 上 proper である．
\end{lemma}

\begin{proof}
Quot, Lemma \ref{quot-lemma-hilb-is-quot} により
$\Hilbfunctor_{X/B} = \Quotfunctor_{\mathcal{O}_X/X/B}$ であることを想起せよ．
したがって，この補題は Lemma \ref{moduli-lemma-quot-proper-over-base} の直ちに得られる帰結である．
\end{proof}

\begin{lemma}
\label{moduli-lemma-hilb-qc-over-base}
$f : X \to B$ を，代数空間の分離的かつ有限表示な射とする．
$\mathcal{L}$ を $X/B$ 上豊富な可逆 $\mathcal{O}_X$-加群とする；
Divisors on Spaces, Definition
\ref{spaces-divisors-definition-relatively-ample} を参照せよ．
$\mathcal{L}$ に関する Hilbert 多項式が $P$ である閉部分スキームを
パラメータ付ける代数空間 $\Hilbfunctor^P_{X/B}$ は，$B$ 上分離的かつ有限表示である．
\end{lemma}

\begin{proof}
Quot, Lemma \ref{quot-lemma-hilb-is-quot} により
$\Hilbfunctor_{X/B} = \Quotfunctor_{\mathcal{O}_X/X/B}$ であることを想起せよ．
したがって，この補題は Lemma \ref{moduli-lemma-quot-qc-over-base} の直ちに得られる帰結である．
\end{proof}
% END EXPANDED MODULE ch108_04.tex
% BEGIN EXPANDED MODULE ch108_05.tex
\section{Picard スタックの性質}
\label{moduli-section-picard-stack}

\noindent
$f : X \to B$ を，平坦，proper，かつ有限表示な代数空間の射とする．
このとき $X/B$ 上の可逆層をパラメータ付けるスタック
$\Picardstack_{X/B}$ は代数的である；Quot, Proposition
\ref{quot-proposition-pic} を参照せよ．

\begin{lemma}
\label{moduli-lemma-pic-diagonal-affine-fp}
$\Picardstack_{X/B}$ の $B$ 上の対角射はアフィンかつ有限表示である．
\end{lemma}

\begin{proof}
Quot, Lemma \ref{quot-lemma-picard-stack-open-in-coh} で，
$\Picardstack_{X/B}$ は $\Cohstack_{X/B}$ の開部分スタックであることを見た．
したがって Lemma \ref{moduli-lemma-coherent-diagonal-affine-fp} から従う．
\end{proof}

\begin{lemma}
\label{moduli-lemma-pic-qs-lfp}
射 $\Picardstack_{X/B} \to B$ は準分離的かつ局所有限表示である．
\end{lemma}

\begin{proof}
Quot, Lemma \ref{quot-lemma-picard-stack-open-in-coh} で，
$\Picardstack_{X/B}$ は $\Cohstack_{X/B}$ の開部分スタックであることを見た．
したがって Lemma \ref{moduli-lemma-coherent-qs-lfp} から従う．
\end{proof}

\begin{lemma}
\label{moduli-lemma-pic-existence-part}
$X \to B$ は proper であることに加えて平滑であると仮定する．
このとき $\Picardstack_{X/B} \to B$ は，付値判定法
(Morphisms of Stacks, Definition
\ref{stacks-morphisms-definition-existence}) の存在部分を満たす．
\end{lemma}

\begin{proof}
基底変換を行えば，これは直ちに次の問題に帰着する：付値環 $R$ とその分数体 $K$，
$R$ 上 proper かつ平滑な代数空間 $X$，および可逆
$\mathcal{O}_{X_K}$-加群 $\mathcal{L}_K$ が与えられたとき，
一般ファイバーが $\mathcal{L}_K$ である可逆 $\mathcal{O}_X$-加群
$\mathcal{L}$ が存在することを示せ．
$X_K$ は Noether 的，分離的，かつ正則であることに注意する
(Morphisms of Spaces, Lemma
\ref{spaces-morphisms-lemma-finite-presentation-noetherian} および
Spaces over Fields, Lemma \ref{spaces-over-fields-lemma-smooth-regular} を用いる)．
したがって，Divisors on Spaces, Lemma
\ref{spaces-divisors-lemma-Noetherian-regular-separated-pic-effective-Cartier} により，
$X_K$ の二つの有効 Cartier 因子 $D_K, D'_K$ に対する
$\mathcal{O}_{X_K}(D_K)$ と $\mathcal{O}_{X_K}(D'_K)$ の Picard 群での差として
$\mathcal{L}_K$ を書ける．最後に，Divisors on Spaces, Lemma
\ref{spaces-divisors-lemma-smooth-over-valuation-ring-effective-Cartier} により，
$D_K$ と $D'_K$ は有効 Cartier 因子 $D, D' \subset X$ の制限であることが分かる．
\end{proof}

\begin{lemma}
\label{moduli-lemma-pic-inertia}
$B$ 上のすべてのスキーム $T$ に対して
$f_{T, *}\mathcal{O}_{X_T} \cong \mathcal{O}_T$ であると仮定する．
このとき $\Picardstack_{X/B}$ の慣性スタックは
$\mathbf{G}_m \times \Picardstack_{X/B}$ に等しい．
\end{lemma}

\begin{proof}
これは Examples of Stacks, Example
\ref{examples-stacks-example-inertia-stack-of-picard} で説明されている．
\end{proof}

\begin{lemma}
\label{moduli-lemma-pic-curves-smooth}
本節の他の仮定に加えて，$f : X \to B$ の相対次元が $\leq 1$ であると仮定する．
このとき $\Picardstack_{X/B} \to B$ は平滑である．
\end{lemma}

\begin{proof}
$\Picardstack_{X/B} \to B$ が局所有限表示であることは既に分かっている；
Lemma \ref{moduli-lemma-pic-qs-lfp} を参照せよ．したがって
$\Picardstack_{X/B} \to B$ が形式的に平滑であることを示せば十分である；
More on Morphisms of Stacks, Lemma
\ref{stacks-more-morphisms-lemma-smooth-formally-smooth} を参照せよ．
基底変換を行えば，これは直ちに次の問題に帰着する：アフィンスキームの一次厚化
$T \subset T'$，proper，平坦，有限表示，かつ相対次元 $\leq 1$ の射
$X' \to T'$，および $X = T \times_{T'} X'$ 上の可逆
$\mathcal{O}_X$-加群 $\mathcal{L}$ が与えられたとき，$X$ への制限が
$\mathcal{L}$ となる可逆 $\mathcal{O}_{X'}$-加群 $\mathcal{L}'$ の存在を示せ．
$T \subset T'$ は一次厚化なので，More on Morphisms of Spaces, Lemma
\ref{spaces-more-morphisms-lemma-base-change-thickening} により
$X \subset X'$ についても同じことが成り立つ．
More on Morphisms of Spaces, Lemma
\ref{spaces-more-morphisms-lemma-picard-group-first-order-thickening} により，
$X$ を $X'$ の中で切り出す準連接イデアルを $\mathcal{I}$ として，
$H^2(X, \mathcal{I}) = 0$ を示せば十分である．
% P12-ERR-0211: Japanese uses a grammatical notation construction for the frozen missing “by”.
構造射を $f : X \to T$ と記す．Cohomology of Spaces, Lemma
\ref{spaces-cohomology-lemma-higher-direct-images-zero-above-dimension-fibre} により，
$p > 1$ に対して $R^pf_*\mathcal{I} = 0$ である．したがって，Cohomology of Spaces, Lemma
\ref{spaces-cohomology-lemma-quasi-coherence-higher-direct-images-application} により
求める消滅を得る（ここで最後に $T$ がアフィンであることを用いる）．
\end{proof}
% END EXPANDED MODULE ch108_05.tex
% BEGIN EXPANDED MODULE ch108_06.tex
\section{Picard 関手の性質}
\label{moduli-section-picard-functor}

\noindent
$f : X \to B$ を平坦，proper，かつ有限表示な代数空間の射とし，さらにすべての
$T/B$ に対して標準写像
$$
\mathcal{O}_T \longrightarrow f_{T, *}\mathcal{O}_{X_T}
$$
が同型であると仮定する．このとき Picard 関手 $\Picardfunctor_{X/B}$ は代数空間である；
Quot, Proposition \ref{quot-proposition-pic-functor} を参照せよ．
% P12-ERR-0212: Japanese renders the intended “close relationship” despite the frozen “closed”.
Picard スタックとの間には密接な関係がある．

\begin{lemma}
\label{moduli-lemma-pic-gerbe-over-pic-functor}
射 $\Picardstack_{X/B} \to \Picardfunctor_{X/B}$ により，Picard スタックは
Picard 関手上の gerbe となる．
\end{lemma}

\begin{proof}
$\Picardstack_{X/B} \to \Picardfunctor_{X/B}$ が gerbe であることの定義は
Morphisms of Stacks, Definition \ref{stacks-morphisms-definition-gerbe} にあり，
そこでさらに Stacks, Definition
\ref{stacks-definition-gerbe-over-stack-in-groupoids} が参照されている．
これを示すため，Stacks, Lemma \ref{stacks-lemma-when-gerbe} の条件 (2)(a) と (2)(b) を確認する．
Quot, Lemma \ref{quot-lemma-pic-over-pic} から直ちに従うが，詳しく説明する．

\medskip\noindent
条件 (2)(a)．$B$ 上のあるスキーム $U$ に対して
$\xi \in \Picardfunctor_{X/B}(U)$ と仮定する．
$\Picardfunctor_{X/B}$ は，$B$ 上のスキームにおける規則
$T \mapsto \Pic(X_T)$ の fppf 層化である
(Quot, Situation \ref{quot-situation-pic})．したがって，ある fppf 被覆
$\{U_i \to U\}$ が存在して，$\xi|_{U_i}$ は $X_{U_i}$ 上のある可逆加群
$\mathcal{L}_i$ に対応する．このとき $(U_i \to B, \mathcal{L}_i)$ は
$U_i$ 上の $\Picardstack_{X/B}$ の対象で，$\xi|_{U_i}$ へ写る．

\medskip\noindent
条件 (2)(b)．$U$ を $B$ 上のスキームとし，$\mathcal{L}, \mathcal{N}$ を
$X_U$ 上の可逆加群で，$\Picardfunctor_{X/B}(U)$ の同じ元へ写るものとする．
このとき，ある fppf 被覆 $\{U_i \to U\}$ が存在して，
$\mathcal{L}|_{X_{U_i}}$ と $\mathcal{N}|_{X_{U_i}}$ は同型である．
したがって求めるとおり，同型
$(U \to B, \mathcal{L})|_{U_i} \to (U \to B, \mathcal{N})|_{U_i}$ を得る．
\end{proof}

\begin{lemma}
\label{moduli-lemma-pic-functor-diagonal-qc-immersion}
$\Picardfunctor_{X/B}$ の $B$ 上の対角射は準コンパクトな immersion である．
\end{lemma}

\begin{proof}
Quot, Lemma \ref{quot-lemma-diagonal-pic} により，対角射は immersion である．
残りとして，対角射が準コンパクトであることを示す．
Lemma \ref{moduli-lemma-pic-diagonal-affine-fp} により $\Picardstack_{X/B}$ の対角射は準コンパクトであり，
Lemma \ref{moduli-lemma-pic-gerbe-over-pic-functor} により
$\Picardstack_{X/B}$ は $\Picardfunctor_{X/B}$ 上の gerbe である．
Morphisms of Stacks, Lemma
\ref{stacks-morphisms-lemma-gerbe-diagonal-quasi-compact} により結論を得る．
\end{proof}

\begin{lemma}
\label{moduli-lemma-pic-functor-qs-lfp}
射 $\Picardfunctor_{X/B} \to B$ は準分離的かつ局所有限表示である．
\end{lemma}

\begin{proof}
$\Picardfunctor_{X/B} \to B$ が準分離的であることを確認するには，
その対角射が準コンパクトであることを示せばよい．これは Lemma
\ref{moduli-lemma-pic-functor-diagonal-qc-immersion} から直ちに従う．
射 $\Picardstack_{X/B} \to \Picardfunctor_{X/B}$ は全射，平坦，かつ局所有限表示である
(Lemma \ref{moduli-lemma-pic-gerbe-over-pic-functor} および Morphisms of Stacks, Lemma
\ref{stacks-morphisms-lemma-gerbe-fppf} による)．したがって Morphisms of Stacks, Lemma
\ref{stacks-morphisms-lemma-flat-finite-presentation-permanence} により，
$\Picardstack_{X/B} \to B$ が局所有限表示であることを示せば十分である．
これは Lemma \ref{moduli-lemma-pic-qs-lfp} から従う．
\end{proof}

\begin{lemma}
\label{moduli-lemma-pic-functor-uniqueness-part}
本節の他の仮定に加えて，$X \to B$ の幾何ファイバーが整であると仮定する．
このとき $\Picardfunctor_{X/B} \to B$ は分離的である．
\end{lemma}

\begin{proof}
$\Picardfunctor_{X/B} \to B$ は準分離的なので，付値判定法の一意性部分を確認すれば十分である；
Morphisms of Spaces, Lemma
\ref{spaces-morphisms-lemma-valuative-criterion-separatedness} を参照せよ．
これは直ちに次の問題に帰着する：
\begin{enumerate}
\item 分数体 $K$ をもつ付値環 $R$，
% P12-ERR-0213: Japanese number marking is neutral; the intended all-geometric-fibres condition is recorded in the ledger.
\item $R$ 上 proper かつ平坦で，幾何ファイバーが整である代数空間 $X$，
\item $a|_{\Spec(K)} = 0$ を満たす元 $a \in \Picardfunctor_{X/R}(R)$
\end{enumerate}
が与えられたとき，$a = 0$ を示さなければならない．
全射かつ平坦な射
$\Picardstack_{X/R} \to \Picardfunctor_{X/R}$ に
Morphisms of Stacks, Lemma
\ref{stacks-morphisms-lemma-lift-valuation-ring-through-flat-morphism} を適用する
（全射かつ平坦であることは Lemma \ref{moduli-lemma-pic-gerbe-over-pic-functor} と
Morphisms of Stacks, Lemma \ref{stacks-morphisms-lemma-gerbe-fppf} による）．
$R$ を拡大に置き換えた後，$a$ は可逆 $\mathcal{O}_X$-加群
$\mathcal{L}$ によって与えられると仮定してよい．
$a|_{\Spec(K)} = 0$ なので，Quot, Lemma
\ref{quot-lemma-flat-geometrically-connected-fibres} により
$\mathcal{L}_K \cong \mathcal{O}_{X_K}$ である．

\medskip\noindent
% P12-ERR-0214: Japanese uses a grammatical notation construction for the frozen missing “by”.
構造射を $f : X \to \Spec(R)$ と記す．
$\eta, 0 \in \Spec(R)$ をそれぞれ一般点と閉点とする．
$\Spec(R)$ 上の perfect 複体
$K = Rf_*\mathcal{L}$ および $M = Rf_*(\mathcal{L}^{\otimes -1})$ を考える；
Derived Categories of Spaces, Lemma
\ref{spaces-perfect-lemma-flat-proper-perfect-direct-image-general} を参照せよ．
$K$ と $M$ に付随する，Derived Categories of Spaces, Lemma
\ref{spaces-perfect-lemma-jump-loci} の函数
$\beta_{K, i}, \beta_{M, i} : \Spec(R) \to \mathbf{Z}$ を考える．
% P12-ERR-0215: Japanese renders the intended conjunction despite the frozen typo “amd”.
$K$ と $M$ の形成は基底変換と可換するので（上で引用した補題を参照），
Spaces over Fields, Lemma
\ref{spaces-over-fields-lemma-proper-geometrically-reduced-global-sections} と
$f$ のファイバーに関する仮定により
$\beta_{K, 0}(\eta) = \beta_{M, 0}(\eta) = 1$ である．
上半連続性により
$\beta_{K, 0}(0) \geq 1$ および $\beta_{M, 0}(0) \geq 1$ を得る．
Spaces over Fields, Lemma
\ref{spaces-over-fields-lemma-characterize-trivial-pic-integral} により，
$\mathcal{L}$ の特殊ファイバー $X_0$ への制限は自明であると結論できる．
これから，上と同様に
% P12-ERR-0216: The frozen unevaluated beta_{M,0} term is retained literally.
$\beta_{K, 0}(0) = \beta_{M, 0} = 1$ を得る．
すると More on Algebra, Lemma
\ref{more-algebra-lemma-lift-pseudo-coherent-from-residue-field} により，
$K$ は次の形の複体によって表現できる：
$$
\ldots \to 0 \to R \to R^{\oplus \beta_{K, 1}(0)} \to
R^{\oplus \beta_{K, 2}(0)} \to \ldots
$$
$\beta_{K, 0}(\eta) = 1$ なので，$R \to R^{\oplus \beta_{K, 1}(0)}$ は零である．
言い換えれば，$D(R)$ において $K = R \oplus \tau_{\geq 1}(K)$ であり，
$\tau_{\geq 1}(K)$ はある $b \in \mathbf{Z}$ に対して $[1, b]$ 内の tor amplitude をもつ．
したがって大域切断 $s \in H^0(X, \mathcal{L})$ で，$X_0$ への制限
$s_0$ が至る所零でないものが存在する（ここでも $K$ の形成が基底変換と可換することを用いる）．
すると $s : \mathcal{O}_X \to \mathcal{L}$ は可逆層の写像であり，
$X_0$ への制限は同型であるから，求めるとおり同型である．
\end{proof}

\begin{lemma}
\label{moduli-lemma-pic-functor-curves-smooth}
本節の他の仮定に加えて，$f : X \to B$ の相対次元が $\leq 1$ であると仮定する．
このとき $\Picardfunctor_{X/B} \to B$ は平滑である．
\end{lemma}

\begin{proof}
Lemma \ref{moduli-lemma-pic-curves-smooth} により
$\Picardstack_{X/B} \to B$ は平滑である．
Lemma \ref{moduli-lemma-pic-gerbe-over-pic-functor} と Morphisms of Stacks, Lemma
\ref{stacks-morphisms-lemma-gerbe-smooth} を組み合わせれば，射
$\Picardstack_{X/B} \to \Picardfunctor_{X/B}$ は全射かつ平滑である．
したがって，$U$ がスキームで $U \to \Picardstack_{X/B}$ が全射かつ平滑ならば，
$U \to \Picardfunctor_{X/B}$ は全射かつ平滑であり，
$U \to B$ も全射かつ平滑である（これらの性質は合成で保たれる）．
ゆえに，例えば Descent on Spaces, Lemma
\ref{spaces-descent-lemma-syntomic-smooth-etale-permanence} により
$\Picardfunctor_{X/B} \to B$ は平滑である．
\end{proof}
% END EXPANDED MODULE ch108_06.tex
% BEGIN EXPANDED MODULE ch108_07.tex
\section{相対射の性質}
\label{moduli-section-relative-morphisms}

\noindent
$B$ を代数空間とする．$X$ と $Y$ を $B$ 上の代数空間で，
$Y \to B$ は平坦，proper，かつ有限表示であり，
$X \to B$ は分離的かつ有限表示であるものとする．
このとき相対射の関手 $\mathit{Mor}_B(Y, X)$ は，
$B$ 上局所有限表示な代数空間である．
Quot, Proposition \ref{quot-proposition-Mor} を参照せよ．

\begin{lemma}
\label{moduli-lemma-Mor-diagonal-closed}
$\mathit{Mor}_B(Y, X) \to B$ の対角射は有限表示な閉埋め込みである．
\end{lemma}

\begin{proof}
Quot, Lemma \ref{quot-lemma-Mor-into-Hilb-open} により，開埋め込み
$\mathit{Mor}_B(Y, X) \to \Hilbfunctor_{Y \times_B X/B}$ が存在する．
したがって Lemma \ref{moduli-lemma-hilb-diagonal-closed} から従う．
\end{proof}

\begin{lemma}
\label{moduli-lemma-Mor-s-lfp}
射 $\mathit{Mor}_B(Y, X) \to B$ は分離的かつ局所有限表示である．
\end{lemma}

\begin{proof}
$\mathit{Mor}_B(Y, X) \to B$ が分離的であることを確認するには，
その対角射が閉埋め込みであることを示せばよい．これは Lemma
\ref{moduli-lemma-Mor-diagonal-closed} により成り立つ．第二の主張は
Quot, Proposition \ref{quot-proposition-Mor} の一部である．
\end{proof}

\begin{lemma}
\label{moduli-lemma-Isom-in-Mor}
$B, X, Y$ は本節の序文におけるものとし，さらに $X \to B$ が
proper であると仮定する．このとき，同型射からなる部分関手
$\mathit{Isom}_B(Y, X) \subset \mathit{Mor}_B(Y, X)$ は開部分空間である．
\end{lemma}

\begin{proof}
More on Morphisms of Spaces, Lemma
\ref{spaces-more-morphisms-lemma-where-isomorphism} から直ちに従う．
\end{proof}

\begin{remark}[数値的不変量]
\label{moduli-remark-Mor-numerical}
$B, X, Y$ は本節の序文におけるものとする．
$I$ を集合とし，各 $i \in I$ に対して
$E_i \in D(\mathcal{O}_{Y \times_B X})$ を perfect とする．
$P : I \to \mathbf{Z}$ を函数とする．Quot, Lemma
\ref{quot-lemma-Mor-into-Hilb-open} により，
$$
\mathit{Mor}_B(Y, X) \subset
\Hilbfunctor_{Y \times_B X/B}
$$
は開部分空間であることを想起せよ．したがって，
$$
\mathit{Mor}^P_B(Y, X) =
\mathit{Mor}_B(Y, X) \cap \Hilbfunctor^P_{Y \times_B X/B}
$$
と定義できる．ここで $\Hilbfunctor^P_{Y \times_B X/B}$ は Remark
\ref{moduli-remark-hilb-numerical} におけるものである．射
$$
\mathit{Mor}^P_B(Y, X) \longrightarrow \mathit{Mor}_B(Y, X)
$$
は平坦な閉埋め込みであり，例えば $I$ が有限であるか，$B$ が局所 Noether 的であるか，
または，ある可逆 $\mathcal{O}_{Y \times_B X}$-加群 $\mathcal{L}$ に対して
$I = \mathbf{Z}$，$E_i = \mathcal{L}^{\otimes i}$ である場合には，開かつ閉な埋め込みである．
最後の場合には，記法 $\mathit{Mor}^{P, \mathcal{L}}_B(Y, X)$ を用いることがある．
\end{remark}

\begin{lemma}
\label{moduli-lemma-Mor-qc-over-base}
$B, X, Y$ は本節の序文におけるものとし，$\mathcal{L}$ を $X/B$ 上豊富，
$\mathcal{N}$ を $Y/B$ 上豊富とする．Divisors on Spaces, Definition
\ref{spaces-divisors-definition-relatively-ample} を参照せよ．
$P$ を数値多項式とする．このとき，
$$
\mathit{Mor}^{P, \mathcal{M}}_B(Y, X) \longrightarrow B
$$
は分離的かつ有限表示である．ここで
$\mathcal{M} = \text{pr}_1^*\mathcal{N}
\otimes_{\mathcal{O}_{Y \times_B X}} \text{pr}_2^*\mathcal{L}$ である．
\end{lemma}

\begin{proof}
Lemma \ref{moduli-lemma-Mor-s-lfp} により，射 $\mathit{Mor}_B(Y, X) \to B$ は
分離的かつ局所有限表示である．したがって，Remark
\ref{moduli-remark-Mor-numerical} の開かつ閉な部分空間
$\mathit{Mor}^{P, \mathcal{M}}_B(Y, X)$ が $B$ 上準コンパクトであることを
示せば十分である．

\medskip\noindent
この問題は $B$ 上 \'etale 局所的である
(Morphisms of Spaces, Lemma \ref{spaces-morphisms-lemma-quasi-compact-local})．
したがって $B$ はアフィンであると仮定してよい．

\medskip\noindent
$B = \Spec(\Lambda)$ と仮定する．$X$ と $Y$ はスキームであり，
$\mathcal{L}$ と $\mathcal{N}$ はそれぞれ $X$ と $Y$ 上の豊富な可逆層であることに
注意する（これは定義から直ちに従う）．$\Lambda = \colim \Lambda_i$ を，
その有限型 $\mathbf{Z}$-部分代数の帰納極限として書く．このとき，ある $i$ と，
$B_i = \Spec(\Lambda_i)$ 上で補題におけるような系
$X_i, Y_i, \mathcal{L}_i, \mathcal{N}_i$ で，$B$ への基底変換が
$X, Y, \mathcal{L}, \mathcal{N}$ を与えるものを見いだせる．これは Limits, Lemmas
\ref{limits-lemma-descend-finite-presentation}（$X_i$, $Y_i$ を見いだすため），
\ref{limits-lemma-descend-invertible-modules}（$\mathcal{L}_i$,
$\mathcal{N}_i$ を見いだすため），
\ref{limits-lemma-descend-separated-finite-presentation}
（$X_i \to B_i$ を分離的にするため），
\ref{limits-lemma-eventually-proper}（$Y_i \to B_i$ を proper にするため），および
\ref{limits-lemma-limit-ample}（$\mathcal{L}_i$, $\mathcal{N}_i$ を豊富にするため）
から従う．
$$
\mathit{Mor}_B(Y, X) = B \times_{B_i} \mathit{Mor}_{B_i}(Y_i, X_i)
$$
であり，さらに
% P12-ERR-0217: The frozen shortened Mor^P notation is retained rather than silently restored to Mor^{P,M}.
$\mathit{Mor}^P_B(Y, X)$ についても同様なので，次の段落で扱う場合に帰着する．

\medskip\noindent
$B$ を Noether アフィンスキームと仮定する．Properties, Lemma
\ref{properties-lemma-ample-on-product} により，$\mathcal{M}$ は豊富である．
Lemma \ref{moduli-lemma-hilb-qc-over-base} により，
$\Hilbfunctor^{P, \mathcal{M}}_{Y \times_B X/B}$ は $B$ 上有限表示であり，
したがって Noether 的である．構成により
$$
\mathit{Mor}^{P, \mathcal{M}}_B(Y, X) =
\mathit{Mor}_B(Y, X) \cap
\Hilbfunctor^{P, \mathcal{M}}_{Y \times_B X/B}
$$
は $\Hilbfunctor^{P, \mathcal{M}}_{Y \times_B X/B}$ の開部分空間であり，
したがって準コンパクトである（Noether 代数空間の開部分空間は準コンパクトであるため）．
\end{proof}
% END EXPANDED MODULE ch108_07.tex
% BEGIN EXPANDED MODULE ch108_08.tex
\section{偏極付き proper スキームのスタックの性質}
\label{moduli-section-polarized}

\noindent
本節では，モジュライスタック
$$
\Polarizedstack \longrightarrow \Spec(\mathbf{Z})
$$
の性質を論じる．スキーム $S$ 上のその切断の圏は，相対的に豊富な可逆層を備えた，
$S$ 上 proper，平坦，かつ有限表示なスキームの圏である．
% P12-ERR-0218: Japanese number marking is neutral for the frozen singular “scheme”.
これは Quot, Theorem \ref{quot-theorem-polarized-algebraic} により代数スタックである．

\begin{lemma}
\label{moduli-lemma-polarized-diagonal-separated-fp}
$\Polarizedstack$ の対角射は分離的かつ有限表示である．
\end{lemma}

\begin{proof}
$\Polarizedstack$ は極限を保つ代数スタックであることを想起せよ；
Quot, Lemma \ref{quot-lemma-polarized-limits} を参照せよ．
Limits of Stacks, Lemma \ref{stacks-limits-lemma-limit-preserving-diagonal} により，
これは
$\Delta : \Polarizedstack \to \Polarizedstack \times \Polarizedstack$
が極限を保つことを意味する．したがって Limits of Stacks, Proposition
\ref{stacks-limits-proposition-characterize-locally-finite-presentation} により，
$\Delta$ は局所有限表示である．

\medskip\noindent
$\Delta$ が分離的であることを示そう．このためには，アフィンスキーム $U$ と，
$U$ 上の $\Polarizedstack$ の二つの対象
$\upsilon = (Y, \mathcal{N})$ および $\chi = (X, \mathcal{L})$ が与えられたとき，
代数空間
$$
\mathit{Isom}_{\Polarizedstack}(\upsilon, \chi)
$$
が分離的であることを示せば十分である．同型 $\upsilon_T \to \chi_T$ に
その台となる同型 $Y_T \to X_T$ を対応させる規則は射
$$
\mathit{Isom}_{\Polarizedstack}(\upsilon, \chi)
\longrightarrow
\mathit{Isom}_U(Y, X)
$$
を定める．Lemmas \ref{moduli-lemma-Mor-s-lfp} および
\ref{moduli-lemma-Isom-in-Mor} で，標的が分離的な代数空間であることを見たので，
この射が分離的であることを示せば十分である．あるスキーム $T/U$ 上の同型
$f : Y_T \to X_T$ が与えられたとき，明らかに
% P12-ERR-0219: Japanese uses a grammatical conditional construction.
$$
\mathit{Isom}_{\Polarizedstack}(\upsilon, \chi)
\times_{\mathit{Isom}_U(Y, X), [f]} T
=
\mathit{Isom}(\mathcal{N}_T, f^*\mathcal{L}_T)
$$
である．ここで $[f] : T \to \mathit{Isom}_U(Y, X)$ は $f$ に対応する
$T$-値点を表し，$\mathit{Isom}(\mathcal{N}_T, f^*\mathcal{L}_T)$ は
Section \ref{moduli-section-hom-isom} で論じた代数空間である．
% P12-ERR-0220: The frozen base U is retained rather than silently changed to T.
この代数空間は $U$ 上アフィンなので，この主張から $\Delta$ が分離的であることが従う．

\medskip\noindent
証明を完了するため，$\Delta$ が準コンパクトであることを示す．
$\Delta$ は代数空間によって表現可能なので，全射かつ平滑な射
$U \to \Polarizedstack \times \Polarizedstack$ による $\Delta$ の基底変換が
準コンパクトであることを確認すれば十分である
% P12-ERR-0221: Japanese normalizes the frozen subject–verb disagreement.
（例えば Properties of Stacks, Lemma
\ref{stacks-properties-lemma-check-property-covering} を参照せよ）．
$U = \coprod U_i$ はアフィン開集合の非交和であると仮定してよい．
$\Polarizedstack$ は極限を保つので（上を参照），
$\Polarizedstack \to \Spec(\mathbf{Z})$ は局所有限表示であり，したがって
$U_i \to \Spec(\mathbf{Z})$ は局所有限表示である
（Limits of Stacks, Proposition
\ref{stacks-limits-proposition-characterize-locally-finite-presentation} および
Morphisms of Stacks, Lemmas
\ref{stacks-morphisms-lemma-composition-finite-presentation} と
\ref{stacks-morphisms-lemma-smooth-locally-finite-presentation}）．
特に $U_i$ は Noether アフィンスキームである．これにより次の段落で扱う場合に帰着する．

\medskip\noindent
この段落では，Noether アフィンスキーム $U$ と，$U$ 上の $\Polarizedstack$ の二つの対象
$\upsilon = (Y, \mathcal{N})$ および $\chi = (X, \mathcal{L})$ が与えられたとき，
代数空間
$$
\mathit{Isom}_{\Polarizedstack}(\upsilon, \chi)
$$
が準コンパクトであることを示す．$U$ の連結成分は開かつ閉なので，
$U$ をそれらで置き換えてよい．したがって $U$ は連結であると仮定する．
$u \in U$ を点とする．$P$ を Hilbert 多項式
$n \mapsto \chi(Y_u, \mathcal{N}_u^{\otimes n})$ とする；
Varieties, Lemma \ref{varieties-lemma-numerical-polynomial-from-euler} を参照せよ．
$U$ は連結であり，函数
$u \mapsto \chi(Y_u, \mathcal{N}_u^{\otimes n})$ は局所定数なので
（Derived Categories of Schemes, Lemma
\ref{perfect-lemma-chi-locally-constant-geometric} を参照），
$U$ のすべての点で同じ Hilbert 多項式を得る．
$Y \times_U X$ 上で
$\mathcal{M} = \text{pr}_1^*\mathcal{N}
\otimes_{\mathcal{O}_{Y \times_U X}} \text{pr}_2^*\mathcal{L}$
とおく．$U$ 上のあるスキーム $T$ に対して
$(f, \varphi) \in \mathit{Isom}_{\Polarizedstack}(\upsilon, \chi)(T)$
が与えられたとき，すべての $t \in T$ に対して
% P12-ERR-0222: Japanese uses a grammatical conditional construction.
$$
\chi(Y_t, (\text{id} \times f)^*\mathcal{M}^{\otimes n}) =
\chi(Y_t,
\mathcal{N}_t^{\otimes n} \otimes_{\mathcal{O}_{Y_t}}
f_t^*\mathcal{L}_t^{\otimes n}) =
\chi(Y_t, \mathcal{N}_t^{\otimes 2n}) = P(2n)
$$
である．中央の等号では，同型
$\varphi : f^*\mathcal{L}_T \to \mathcal{N}_T$ を用いた．
$P'(t) = P(2t)$ とおくと，上の射
$$
\mathit{Isom}_{\Polarizedstack}(\upsilon, \chi)
\longrightarrow
\mathit{Isom}_U(Y, X)
$$
の像は共通部分
$$
\mathit{Isom}_U(Y, X) \cap \mathit{Mor}^{P', \mathcal{M}}_U(Y, X)
$$
に含まれることが分かる．この共通部分は $\mathit{Mor}_U(Y, X)$ の
開部分空間どうしの共通部分である
（Lemma \ref{moduli-lemma-Isom-in-Mor} と Remark \ref{moduli-remark-Mor-numerical} を参照）．
Lemma \ref{moduli-lemma-Mor-qc-over-base} により
$\mathit{Mor}^{P', \mathcal{M}}_U(Y, X)$ は $U$ 上有限表示なので，
Noether 代数空間である．したがってこの共通部分も Noether 代数空間である．射
$$
\mathit{Isom}_{\Polarizedstack}(\upsilon, \chi)
\longrightarrow
\mathit{Isom}_U(Y, X) \cap \mathit{Mor}^{P', \mathcal{M}}_U(Y, X)
$$
はアフィンなので（上を参照），結論を得る．
\end{proof}

\begin{lemma}
\label{moduli-lemma-polarized-qs-lfp}
射 $\Polarizedstack \to \Spec(\mathbf{Z})$ は準分離的かつ局所有限表示である．
\end{lemma}

\begin{proof}
$\Polarizedstack \to \Spec(\mathbf{Z})$ が準分離的であることを確認するには，
その対角射が準コンパクトかつ準分離的であることを示さなければならない．
これは Lemma \ref{moduli-lemma-polarized-diagonal-separated-fp} から直ちに従う．
$\Polarizedstack \to \Spec(\mathbf{Z})$ が局所有限表示であることを示すには，
Limits of Stacks, Proposition
\ref{stacks-limits-proposition-characterize-locally-finite-presentation} により，
$\Polarizedstack$ が極限を保つことを示せば十分である．
これは Quot, Lemma \ref{quot-lemma-polarized-limits} である．
\end{proof}

\begin{lemma}
\label{moduli-lemma-bounded-polarized}
$n \geq 1$ を整数とし，$P$ を数値多項式とする．
$$
T \subset |\Polarizedstack|
$$
を次の性質をもつ部分集合とする：すべての $\xi \in T$ に対して，体 $k$ と，
$\xi$ を表現する $k$ 上の $\Polarizedstack$ の対象 $(X, \mathcal{L})$ で，
次を満たすものが存在する：
\begin{enumerate}
\item $X$ 上の $\mathcal{L}$ の Hilbert 多項式は $P$ であり，
\item 閉埋め込み $i : X \to \mathbf{P}^n_k$ で
$i^*\mathcal{O}_{\mathbf{P}^n}(1) \cong \mathcal{L}$ を満たすものが存在する．
\end{enumerate}
このとき $T$ は Noether 位相空間であり，特に準コンパクトである．
\end{lemma}

\begin{proof}
$|\Polarizedstack|$ は局所 Noether 位相空間であることに注意する；
Morphisms of Stacks, Lemma
\ref{stacks-morphisms-lemma-Noetherian-topology} を参照せよ
（ここではさらに，$\Spec(\mathbf{Z})$ が Noether であること，したがって
Lemma \ref{moduli-lemma-polarized-qs-lfp} と Morphisms of Stacks, Lemma
\ref{stacks-morphisms-lemma-locally-finite-type-locally-noetherian} により
$\Polarizedstack$ が局所 Noether 代数スタックであることを用いる）．
したがって $|\Polarizedstack|$ の任意の準コンパクト部分集合は Noether 位相空間であり，
その任意の部分集合も Noether である；Topology, Lemmas
\ref{topology-lemma-finite-union-Noetherian} と
\ref{topology-lemma-Noetherian} を参照せよ．
ゆえに，$T$ を含む準コンパクト部分集合を見いだせばよい．
% P12-ERR-0223: Japanese normalizes the frozen extraneous article.

\medskip\noindent
Lemma \ref{moduli-lemma-hilb-proper-over-base} により，代数空間
$$
H =
\Hilbfunctor^{P, \mathcal{O}(1)}_{\mathbf{P}^n_\mathbf{Z}/\Spec(\mathbf{Z})}
$$
は $\Spec(\mathbf{Z})$ 上 proper である．Quot, Lemma
\ref{quot-lemma-extend-hilb-to-spaces}
% P12-ERR-0224: The frozen unresolved editorial placeholder is translated and retained.
\footnote{後で（将来の参照をここに挿入），$H$ がスキームであり，したがってこの補題と
Quot, Lemma \ref{quot-lemma-extend-polarized-to-spaces} を用いる必要がないことを見る．}
により，$H$ の恒等射は閉部分空間
$$
Z \subset \mathbf{P}^n_H
$$
に対応する．これは $H$ 上 proper，平坦，かつ有限表示であり，
制限 $\mathcal{N} = \mathcal{O}(1)|_Z$ は $Z/H$ 上相対的に豊富で，
$Z \to H$ のファイバー上で Hilbert 多項式 $P$ をもつ．
特に，対 $(Z \to H, \mathcal{N})$ は射
$$
H \longrightarrow \Polarizedstack
$$
を定め，これはスキームの射 $U \to H$ を族
$(Z_U \to U, \mathcal{N}_U)$ の分類射へ送る；
Quot, Lemma \ref{quot-lemma-extend-polarized-to-spaces} を参照せよ．
% P12-ERR-0225: The frozen closing parenthesis inside the math span is retained literally.
$H$ は Noether 代数空間なので（$\mathbf{Z})$ 上 proper であるため），
$|H|$ は Noether であり，したがって準コンパクトである．写像
$$
|H| \longrightarrow |\Polarizedstack|
$$
は連続なので，その像は準コンパクトである．
したがって $T$ が $|H| \to |\Polarizedstack|$ の像に含まれることを示せば十分である．
しかし，仮定 (1) と (2) はまさにこのことを表している：閉埋め込み
$i : X \to \mathbf{P}^n_k$ で
$i^*\mathcal{O}_{\mathbf{P}^n}(1) \cong \mathcal{L}$ を満たすものを選べば，
$H$ のモジュライ解釈により $H$ の $k$-値点を得る．
% P12-ERR-0226: Japanese normalizes the frozen missing connective construction.
これで補題の証明が完了する．
\end{proof}
% END EXPANDED MODULE ch108_08.tex
% BEGIN EXPANDED MODULE ch108_09.tex
\section{proper 射上の複体のモジュライの性質}
\label{moduli-section-complexes}

\noindent
$f : X \to B$ を proper，平坦，かつ有限表示な代数空間の射とする．
このとき，負次数の自己 Ext が消える相対 perfect 複体をパラメータ付けるスタック
$\Complexesstack_{X/B}$ は代数的である．
Quot, Theorem \ref{quot-theorem-complexes-algebraic} を参照せよ．

\begin{lemma}
\label{moduli-lemma-complexes-diagonal-affine-fp}
$\Complexesstack_{X/B}$ の $B$ 上の対角射はアフィンかつ有限表示である．
\end{lemma}

\begin{proof}
対角射が代数空間によって表現可能であることは Quot, Lemma
\ref{quot-lemma-complexes-diagonal} で示された．その証明から，
次を示せばよいことが分かる：$B$ 上のスキーム $T$ と，
$\Complexesstack_{X/B}$ の $T$ 上のファイバー圏の対象
$(T, E)$ および $(T, E')$ を与える
$E, E' \in D(\mathcal{O}_{X_T})$ が与えられたとき，
$\mathit{Isom}(E, E') \to T$ はアフィンかつ有限表示である．
ここで $\mathit{Isom}(E, E')$ は関手
$$
(\Sch/T)^{opp} \to \textit{Sets},\quad
T' \mapsto \{\varphi : E_{T'} \to E'_{T'}
\text{ isomorphism in }D(\mathcal{O}_{X_{T'}})\}
$$
であり，$E_{T'}$ と $E'_{T'}$ は $E$ と $E'$ の $X_{T'}$ への導来引き戻しである．
規則
$$
(\Sch/T)^{opp} \to \textit{Sets},\quad
% P12-ERR-0227: The frozen E_T and E'_T subscripts are retained rather than silently changed to T'.
T' \mapsto \Hom_{\mathcal{O}_{X_{T'}}}(E_T, E'_T)
$$
によって定義される関手 $H = \SheafHom(E, E')$ を考える．
Quot, Lemma \ref{quot-lemma-complexes-open-neg-exts-vanishing} により，
これは $T$ 上アフィンかつ有限表示な代数空間である．
$H' = \SheafHom(E', E)$，$I = \SheafHom(E, E)$，および
$I' = \SheafHom(E', E')$ についても同じことが成り立つ．したがって，
% P12-ERR-0228: The frozen fibre-product target and composition order are retained literally.
$$
\mathit{Isom}(E, E') = (H' \times_T H) \times_{c, I \times_T I', \sigma} T
$$
であることが分かる．ここで
$c(\varphi', \varphi) = (\varphi \circ \varphi', \varphi' \circ \varphi)$
かつ $\sigma = (\text{id}, \text{id})$ である
（Quot, Proposition \ref{quot-proposition-isom} の証明と比較せよ）．
したがって $\mathit{Isom}(E, E')$ は，$T$ 上アフィンなスキームの
ファイバー積として $T$ 上アフィンである．同様に，
$\mathit{Isom}(E, E')$ は $T$ 上有限表示である．
\end{proof}

\begin{lemma}
\label{moduli-lemma-complexes-qs-lfp}
射 $\Complexesstack_{X/B} \to B$ は準分離的かつ局所有限表示である．
\end{lemma}

\begin{proof}
$\Complexesstack_{X/B} \to B$ が準分離的であることを確認するには，
その対角射が準コンパクトかつ準分離的であることを示さなければならない．
これは Lemma \ref{moduli-lemma-complexes-diagonal-affine-fp} から直ちに従う．
$\Complexesstack_{X/B} \to B$ が局所有限表示であることを示すには，
Limits of Stacks, Proposition
\ref{stacks-limits-proposition-characterize-locally-finite-presentation} により，
$\Complexesstack_{X/B} \to B$ が極限を保つことを示さなければならない．
これは Quot, Lemma \ref{quot-lemma-complexes-limits} から従う
（細部を一つ省略した）．
\end{proof}


% END EXPANDED MODULE ch108_09.tex
